\documentclass[12pt, oneside, openany]{src/thesis}

\begin{document}

% ==================== TRANG ĐẦU (KHÔNG ĐÁNH SỐ TRANG) ==================== %
% Không đánh số trang: Thông tin hội đồng, Lời cảm ơn, Mục lục, Danh mục hình, bảng, viết tắt
\pagestyle{empty}
\assignpagestyle{\chapter}{empty}

\newgeometry{a4paper, top=2.3cm, bottom=2.3cm, left=3.5cm, right=2.0cm}
\begin{titlepage}
    % Khung viền đôi chuẩn đồ án tốt nghiệp UIT (khổ giấy A4)
    \begin{tikzpicture}[remember picture, overlay]
        % Khung ngoài (nét mảnh 0.75pt)
        \draw[line width=0.75pt] 
            ($(current page.north west) + (2.55cm, -1.75cm)$) 
            rectangle 
            ($(current page.south east) + (-1.05cm, 1.75cm)$);
    \end{tikzpicture}

    % Khối tiêu đề trường và khoa sát phía trên khung viền
    \begin{center}
        {\fontsize{15pt}{16pt}\selectfont \textbf{ĐẠI HỌC QUỐC GIA TP. HỒ CHÍ MINH \\ TRƯỜNG ĐẠI HỌC CÔNG NGHỆ THÔNG TIN \\ KHOA KHOA HỌC VÀ KỸ THUẬT THÔNG TIN}\par}
    \end{center}

    % Khoảng cách 1: đến tên sinh viên
    \vfill
    \begin{center}
        {\fontsize{15pt}{16pt}\selectfont \textbf{LÊ CHÂU ĐĂNG KHOA}\par}
    \end{center}

    % Khoảng cách 2: đến tên đề tài
    \vfill
    \begin{center}
        {\fontsize{15pt}{20pt}\selectfont \textbf{ĐỒ ÁN TỐT NGHIỆP}\par}
        
        \vspace{5pt}
        \begin{minipage}{0.9\linewidth}
            \centering
            {\fontsize{16pt}{22pt}\selectfont \textbf{NGHIÊN CỨU PHƯƠNG PHÁP RAG ĐA~PHƯƠNG~THỨC CHO TRỢ~LÝ~ẢO THÔNG~MINH HỖ~TRỢ DẠY~VÀ~HỌC MÔN~KHOA~HỌC~TỰ~NHIÊN CẤP~TRUNG~HỌC~CƠ~SỞ}\par}
            
            \vspace{8pt}
            {\fontsize{13pt}{18pt}\selectfont \textbf{RESEARCH ON A MULTIMODAL RETRIEVAL-AUGMENTED~GENERATION METHOD FOR AN INTELLIGENT~VIRTUAL~ASSISTANT SUPPORTING NATURAL~SCIENCE TEACHING~AND~LEARNING IN~SECONDARY~SCHOOL.}\par}
        \end{minipage}
    \end{center}

    % Khoảng cách 3: đến bậc đào tạo/ngành
    \vfill
    \begin{center}
        {\fontsize{15pt}{18pt}\selectfont \textbf{CỬ NHÂN NGÀNH CÔNG NGHỆ THÔNG TIN}\par}
    \end{center}

    % Khoảng cách 4: đến địa điểm và năm (sát phía dưới khung viền)
    \vfill
    \begin{center}
        {\fontsize{13pt}{17pt}\selectfont \textbf{TP. HỒ CHÍ MINH, 2026}\par}
    \end{center}
\end{titlepage}
\restoregeometry

\input{src/covers/second_title_page}
\chapter*{THÔNG TIN HỘI ĐỒNG CHẤM ĐỒ ÁN TỐT NGHIỆP}

Hội đồng chấm đồ án tốt nghiệp, thành lập theo Quyết định số . . . . . . . . . . . ngày  . . . . . . . . . . . . . . . của Hiệu trưởng Trường Đại học Công nghệ Thông tin.

\chapter*{LỜI CẢM ƠN}
\addcontentsline{toc}{chapter}{\textbf{LỜI CẢM ƠN}}

Để hoàn thành tốt đồ án tốt nghiệp này, em đã nhận được sự quan tâm, giúp đỡ, động viên và hướng dẫn tận tình từ Quý thầy cô, Nhà trường, Trung tâm Phát triển Công nghệ Thông tin (CITD), gia đình và bạn bè.

Trước hết, em xin bày tỏ lòng biết ơn sâu sắc nhất đến thầy \textbf{ThS. Nguyễn Hữu Quyền}, người đã tận tâm định hướng phương pháp nghiên cứu, chỉ bảo tận tình và truyền đạt cho em những kiến thức chuyên môn quý báu trong suốt quá trình thực hiện đề tài. Sự nhiệt huyết và những lời khuyên học thuật sâu sắc của thầy chính là kim chỉ nam giúp em vượt qua những thách thức kỹ thuật và hoàn thiện báo cáo đồ án.

Em xin trân trọng gửi lời cảm ơn đến Ban Giám hiệu Trường Đại học Công nghệ Thông tin --- ĐHQG-HCM, cùng Ban Giám đốc Trung tâm Phát triển Công nghệ Thông tin (CITD) đã tạo mọi điều kiện học tập, nghiên cứu thuận lợi nhất cho em trong suốt khoá học. Em cũng xin gửi lời cảm ơn chân thành đến toàn thể Quý thầy cô đã trực tiếp giảng dạy, trang bị cho em nền tảng kiến thức khoa học máy tính vững chắc. Đó là hành trang vô giá giúp em có đủ năng lực tiếp cận, nghiên cứu và ứng dụng các công nghệ tiên tiến như Trí tuệ nhân tạo, Mô hình ngôn ngữ lớn (LLM) và hệ thống RAG vào giải quyết các bài toán thực tiễn.

Cuối cùng, em xin gửi lời tri ân sâu sắc nhất đến gia đình và bạn bè, những người luôn là điểm tựa vững chắc, không ngừng động viên, sẻ chia và tạo mọi điều kiện tốt nhất để em vững tâm học tập và nghiên cứu.

Dù đã dành nhiều tâm huyết và nỗ lực để hoàn thành đồ án một cách chỉn chu nhất, song do giới hạn về mặt thời gian cũng như kinh nghiệm nghiên cứu, báo cáo khó tránh khỏi những thiếu sót nhất định. Em rất mong nhận được những nhận xét, góp ý quý báu từ Quý thầy cô trong Hội đồng để báo cáo và hệ thống ngày càng hoàn thiện hơn nữa.

Em xin chân thành cảm ơn!
\vspace{1.5cm}

\begin{flushright}
\textbf{Học viên thực hiện} \\
\textit{Lê Châu Đăng Khoa}
\end{flushright}


\renewcommand{\contentsname}{MỤC LỤC}
\tableofcontents
\newpage

\renewcommand{\figurename}{Hình}
\renewcommand{\listfigurename}{DANH SÁCH HÌNH ẢNH}
\phantomsection
\addchapternopage{\listfigurename}
\listoffigures
\newpage

\renewcommand{\tablename}{Bảng}
\renewcommand{\listtablename}{DANH SÁCH BẢNG BIỂU}
\phantomsection
\addchapternopage{\listtablename}
\listoftables
\newpage

\chapter*{DANH MỤC TỪ VIẾT TẮT}
\addchapternopage{DANH MỤC TỪ VIẾT TẮT}

\begin{longtable}{|l|l|p{7.5cm}|}
    \hline
    \textbf{Từ viết tắt} & \textbf{Từ tiếng Anh} & \textbf{Từ tiếng Việt / Giải nghĩa} \\
    \hline
    \endfirsthead
    \hline
    \textbf{Từ viết tắt} & \textbf{Từ tiếng Anh} & \textbf{Từ tiếng Việt / Giải nghĩa} \\
    \hline
    \endhead
    \hline
    \endfoot
    AI & Artificial Intelligence & Trí tuệ nhân tạo \\
    \hline
    LLM & Large Language Model & Mô hình ngôn ngữ lớn \\
    \hline
    RAG & Retrieval-Augmented Generation & Sinh văn bản tăng cường truy xuất \\
    \hline
    NLP & Natural Language Processing & Xử lý ngôn ngữ tự nhiên \\
    \hline
    VQA & Visual Question Answering & Hỏi đáp trực quan \\
    \hline
    OCR & Optical Character Recognition & Nhận dạng ký tự quang học \\
    \hline
    THCS & - & Trung học cơ sở \\
    \hline
    SGK & - & Sách giáo khoa \\
    \hline
    KHTN & - & Khoa học tự nhiên \\
    \hline
    GPU & Graphics Processing Unit & Bộ xử lý đồ họa \\
    \hline
    API & Application Programming Interface & Giao diện lập trình ứng dụng \\
    \hline
    CPU & Central Processing Unit & Bộ xử lý trung tâm \\
    \hline
    ETL & Extract -- Transform -- Load & Quy trình trích xuất, biến đổi và nạp dữ liệu \\
    \hline
    BM25 & Best Matching 25 & Hàm xếp hạng truy xuất theo thống kê từ khoá \\
    \hline
    MRR & Mean Reciprocal Rank & Trung bình nghịch đảo thứ hạng (điểm rank) \\
    \hline
    RRF & Reciprocal Rank Fusion & Hợp nhất theo thứ hạng nghịch đảo \\
    \hline
    CLIP & Contrastive Language--Image Pre-training & Mô hình liên kết ngôn ngữ -- thị giác \\
    \hline
    IR & Information Retrieval & Truy xuất thông tin \\
    \hline
\end{longtable}

\chapter*{DANH MỤC THUẬT NGỮ}
\addchapternopage{DANH MỤC THUẬT NGỮ}

\begin{longtable}{|l|p{9.5cm}|}
    \hline
    \textbf{Thuật ngữ tiếng Anh} & \textbf{Thuật ngữ tiếng Việt tương đương / Định nghĩa} \\
    \hline
    \endfirsthead
    \hline
    \textbf{Thuật ngữ tiếng Anh} & \textbf{Thuật ngữ tiếng Việt tương đương / Định nghĩa} \\
    \hline
    \endhead
    \hline
    \endfoot
    Embedding & Nhúng (Quá trình biểu diễn thực thể thành vector số học liên tục) \\
    \hline
    Vector Database & Cơ sở dữ liệu Vector (Hệ thống lưu trữ và tìm kiếm vector tương đồng) \\
    \hline
    Tech stack & Hệ thống nền tảng công nghệ / Bộ công nghệ triển khai (Tập hợp các thư viện, framework và hạ tầng) \\
    \hline
    Prompt & Lời nhắc (Prompt) (Ngữ cảnh, câu hỏi và chỉ thị đầu vào cho mô hình ngôn ngữ) \\
    \hline
    Hallucination & Hiện tượng ảo giác (Mô hình sinh thông tin sai lệch thiếu căn cứ) \\
    \hline
    Chunking & Phân đoạn / Phân mảnh (Chia nhỏ văn bản thành các phân đoạn độc lập) \\
    \hline
    Information Retrieval & Truy xuất thông tin (IR) \\
    \hline
    Fine-tuning & Tinh chỉnh mô hình (Huấn luyện chuyên sâu trên tập dữ liệu đặc thù) \\
    \hline
    Generative AI & Trí tuệ nhân tạo tạo sinh \\
    \hline
    Hybrid retrieval & Truy xuất lai (Kết hợp đồng thời kênh từ khoá và kênh ngữ nghĩa) \\
    \hline
    Dense / Sparse retrieval & Truy xuất ngữ nghĩa theo vector đặc / Truy xuất từ khóa theo vector thưa \\
    \hline
    Re-ranking & Tái xếp hạng (Đánh giá lại điểm số và sắp xếp tập ứng viên sơ tuyển) \\
    \hline
    Cross-encoder & Bộ mã hoá chéo / Mô hình tương tác chéo (Transformer xử lý đồng thời cặp truy vấn -- tài liệu) \\
    \hline
    Bi-encoder & Kiến trúc hai nhánh mã hóa độc lập / Bi-encoder (Mô hình mã hoá truy vấn và tài liệu tách biệt) \\
    \hline
    Relevance gate & Cổng lọc độ liên quan (Cơ chế lọc ứng viên dựa trên ngưỡng khoảng cách ngữ nghĩa) \\
    \hline
    Ground truth & Dữ liệu chuẩn đối sánh / Nhãn chuẩn xác thực \\
    \hline
    Ablation study & Thực nghiệm bóc tách thành phần (Phân tích đóng góp độc lập của từng cấu phần) \\
    \hline
    LLM-as-a-judge & Phương pháp đánh giá chất lượng bằng mô hình ngôn ngữ lớn độc lập \\
    \hline
    Human-in-the-loop & Quy trình phối hợp chuyên gia rà soát (Chuyên gia trực tiếp thẩm định và kiểm định) \\
    \hline
    Layout segmentation & Phân vùng bố cục trang tài liệu \\
    \hline
    Multimodal RAG & Sinh văn bản tăng cường truy xuất đa phương thức (Tích hợp đồng thời học liệu văn bản và hình ảnh) \\
    \hline
    Late interaction & Cơ chế tương tác trễ (ColBERT tính tương đồng ở mức token sau khi mã hóa ngữ cảnh) \\
    \hline
    Zero-shot retrieval & Truy xuất không cần tinh chỉnh mẫu (Khả năng tìm kiếm tổng quát trên tập dữ liệu mới) \\
    \hline
    Hit@K / Recall@K & Tỷ lệ trúng tại K / Độ phủ tại K (Tỷ lệ câu hỏi tìm thấy tài liệu liên quan trong top K kết quả) \\
    \hline
    Out-of-domain & Ngoài phạm vi kho tri thức (Các câu hỏi không thuộc nội dung sách giáo khoa mục tiêu) \\
    \hline
    Formula degradation & Hiện tượng suy biến công thức (Lỗi phân rã ký hiệu toán -- lý -- hóa và chỉ số dưới qua OCR) \\
    \hline
    Line stitching & Ghép dòng lai phân cấp (Quy trình hợp nhất có điều kiện giữa văn bản OCR và công thức chuyên sâu) \\
    \hline
\end{longtable}

% =================== BẮT ĐẦU ĐÁNH SỐ TRANG Ả-RẬP TỪ TÓM TẮT =================== %
% Đánh số trang ở giữa bên dưới, theo hệ thống số Ả-rập
\mainmatter
\pagenumbering{arabic}
\setcounter{page}{1}

\pagestyle{fancy}
\fancyhf{}
\rhead{}
\lhead{\itshape\small\leftmark}
\cfoot{\thepage}
\renewcommand{\headrulewidth}{0.4pt}

% Kiểu trang plain (trang đầu của mỗi chapter) cũng đánh số ở giữa bên dưới
\fancypagestyle{plain}{%
  \fancyhf{}%
  \cfoot{\thepage}%
  \renewcommand{\headrulewidth}{0pt}%
}
\assignpagestyle{\chapter}{plain}

\chapter*{TÓM TẮT}
\addcontentsline{toc}{chapter}{\textbf{TÓM TẮT}}

Sự phát triển mạnh mẽ của Trí tuệ nhân tạo, đặc biệt là các Mô hình ngôn ngữ lớn (Large Language Models --- LLMs), đã mở ra nhiều tiềm năng ứng dụng đột phá trong lĩnh vực giáo dục. Dẫu vậy, việc tích hợp trực tiếp LLM vào giáo dục phổ thông --- tiêu biểu là môn Khoa học tự nhiên cấp Trung học cơ sở, môn học tích hợp ba phân môn Vật lí, Hoá học và Sinh học --- vẫn vấp phải hai rào cản căn bản: hiện tượng sinh thông tin sai lệch (ảo giác --- hallucination) và sự thiếu hụt cơ chế truy xuất minh bạch từ nguồn sách giáo khoa chuẩn mực trong nước.

Đồ án tốt nghiệp này nghiên cứu và xây dựng một hệ thống trợ lý ảo ứng dụng kỹ thuật sinh văn bản tăng cường truy xuất đa phương thức (Multi-modal Retrieval-Augmented Generation --- RAG \cite{lewis2020retrieval}) dành riêng cho môn Khoa học tự nhiên cấp Trung học cơ sở. Kho tri thức của hệ thống được số hoá toàn diện từ \textbf{12 cuốn sách giáo khoa} lớp 6, 7, 8 và 9 thuộc ba bộ sách hiện hành (Kết nối tri thức với cuộc sống, Chân trời sáng tạo, Cánh Diều): \textbf{2.387/2.399 trang} nội dung được lập chỉ mục thành \textbf{16.515 đoạn văn bản} (sử dụng mô hình nhúng \texttt{BAAI/bge-m3} \cite{chen2024bgem3}, $1024$ chiều) và \textbf{3.881 khung hình} (sử dụng \texttt{clip-vit-base-patch16} \cite{radford2021clip}, $512$ chiều) trên hệ quản trị cơ sở dữ liệu vector ChromaDB, song song với một chỉ mục thưa BM25 \cite{robertson2009bm25} đồng bộ định danh đoạn. Mô hình sinh ngôn ngữ được lựa chọn là Qwen2.5-3B-Instruct \cite{qwen25_2024}. Toàn bộ quy trình trích xuất và lập chỉ mục dữ liệu được thiết kế tối ưu để vận hành hiệu quả trực tiếp trên \textbf{CPU} của máy trạm thông thường mà không phụ thuộc vào phần cứng GPU đắt đỏ.

Đóng góp cốt lõi về mặt phương pháp luận của đồ án tập trung vào \textbf{pha truy xuất}. Hệ thống kết hợp chặt chẽ giữa cơ chế định tuyến theo ý định truy vấn, truy xuất lai (kết hợp BM25 và tìm kiếm ngữ nghĩa) cùng bước tái xếp hạng bằng mô hình cross-encoder \texttt{BAAI/bge-reranker-v2-m3}, với các công tắc độc lập cho phép đánh giá chính xác đóng góp của từng thành phần. Bộ dữ liệu kiểm thử được thiết lập thông qua phương pháp lấy mẫu \textbf{ngẫu nhiên trên toàn bộ tập ngữ liệu}, bao gồm \textbf{240 câu hỏi}: $158$ câu hỏi văn bản, $52$ câu hỏi hình ảnh và $30$ câu hỏi \textbf{ngoài phạm vi} môn học nhằm kiểm chứng trực tiếp khả năng nhận diện ngoại vi và từ chối trả lời, tránh hiện tượng phỏng đoán sai lệch. Ở cấu hình đề xuất (truy xuất lai kết hợp tái xếp hạng), hệ thống đạt chỉ số \textbf{MRR $0{,}7789$} trên $210$ câu hỏi có nhãn nguồn, vượt trội so với các cấu hình đối chứng (BM25 thuần đạt MRR $0{,}6636$; ngữ nghĩa thuần đạt MRR $0{,}5664$; phương pháp lai chưa qua tái xếp hạng đạt MRR $0{,}6947$) trên hầu hết các chỉ số, ngoại trừ một hiện tượng nghịch đảo đáng lưu ý: chỉ số $R@10$ của cấu hình đề xuất ($0{,}2458$) thấp hơn so với cấu hình phương pháp lai chưa tái xếp hạng ($0{,}2643$). Ngoài ra, kết quả thực nghiệm cũng xác nhận rằng cơ chế lọc theo khoảng cách tương đối làm suy giảm hiệu năng của truy xuất lai ở quy mô triển khai thực tế (MRR giảm từ $0{,}7779$ xuống $0{,}7157$), trong khi bước tái xếp hạng luôn mang lại cải thiện tích cực trên cả ba kênh truy xuất.

Chất lượng câu trả lời sinh ra được đánh giá độc lập theo phương pháp LLM-as-a-judge \cite{zheng2023llmjudge} thông qua một tập hợp bốn mô hình ngôn ngữ xoay vòng, phân tách chi tiết theo từng loại câu hỏi: nhóm câu hỏi ngoài phạm vi đạt kết quả gần như tuyệt đối (Tính đúng $4{,}87/5$, Độ trung thực và Độ liên quan $5{,}00/5$); nhóm văn bản đạt chất lượng tốt (Tính đúng $4{,}19/5$, Độ trung thực $4{,}56/5$, Độ liên quan $4{,}59/5$); nhóm câu hỏi dựa trên hình ảnh đạt mức thấp hơn (Tính đúng $3{,}08/5$). Hệ thống đảm bảo tính minh bạch học thuật khi thông tin trích dẫn trang sách được thiết lập mang tính \textit{xác định} từ siêu dữ liệu của đoạn văn bản truy xuất, triệt tiêu hoàn toàn khả năng mô hình tự sinh sai lệch số trang. Báo cáo cũng thẳng thắn phân tích các giới hạn kỹ thuật: giải pháp kiến trúc lai giữa Tesseract và mô hình thị giác--ngôn ngữ dành cho công thức Vật lí, Hoá học đã hoàn thành xử lý trên toàn bộ ngữ liệu ($4\,293$ lượt ghép dòng thành công), song mức độ cải thiện của chỉ số dưới vẫn cần được tiếp tục lượng hoá qua các đợt đo đối sánh chuyên sâu; đồng thời thực nghiệm đối chiếu đa phương thức trên $52$ câu hỏi hình ảnh cho thấy ngữ cảnh đa phương thức chưa mang lại cải thiện rõ rệt (Tính đúng $\Delta{-0{,}077}$, Độ trung thực và Độ liên quan $\Delta{-0{,}038}$).

% \vspace{1em}
% \noindent\textbf{Từ khóa:} Trí tuệ nhân tạo, Mô hình ngôn ngữ lớn (LLM), Sinh văn bản tăng cường truy xuất (RAG), RAG đa phương thức, Truy xuất lai, Trợ lý ảo, Khoa học tự nhiên trung học cơ sở, Qwen2.5, Vector Database.


\chapter{TỔNG QUAN ĐỀ TÀI}
\label{chapter:tong_quan_de_tai}

\section{Bối cảnh và tính cấp thiết}
Trong tiến trình chuyển đổi số sâu rộng của ngành giáo dục, việc nghiên cứu và ứng dụng trí tuệ nhân tạo vào hoạt động dạy và học đang trở thành một xu thế tất yếu, mở ra nhiều cơ hội nâng cao chất lượng sư phạm~\cite{zawacki2019systematic, kasneci2023chatgpt, holmes2019ai, luckin2016intelligence}. Sự bùng nổ của các mô hình ngôn ngữ lớn đã đánh dấu bước phát triển vượt bậc về khả năng thấu hiểu ngữ cảnh và tạo lập ngôn ngữ tự nhiên~\cite{brown2020gpt3, zhao2023llmsurvey}. Dựa trên nền tảng này, các hệ thống trợ lý học tập có thể hỗ trợ người học tra cứu kiến thức, giải thích cặn kẽ các vấn đề khoa học phức tạp và tương tác linh hoạt theo lộ trình cá nhân hóa~\cite{kasneci2023chatgpt}. Đối với môn Khoa học tự nhiên trong Chương trình Giáo dục phổ thông cấp Trung học cơ sở, tiềm năng ứng dụng của các công nghệ này càng thể hiện rõ nét. Là môn học tích hợp liên môn giữa ba phân môn Vật lí, Hoá học và Sinh học, chương trình đòi hỏi học sinh phải tiếp thu đồng thời nhiều dạng tri thức phong phú: từ các cơ chế và quá trình biến đổi sinh học, các định luật vật lí gắn liền với công thức tính toán định lượng, cho đến bản chất các phản ứng hoá học và phương trình biến đổi chất.

Tuy vậy, việc đưa trực tiếp các mô hình ngôn ngữ lớn thương mại vào phục vụ giáo dục phổ thông hiện gặp phải nhiều rào cản kỹ thuật căn bản~\cite{kasneci2023chatgpt}. Trước hết, do được huấn luyện trên nguồn ngữ liệu mở khổng lồ từ Internet, các mô hình ngôn ngữ lớn rất dễ nảy sinh hiện tượng tạo lập thông tin sai lệch, ngụy tạo tri thức nhưng lại được diễn đạt với văn phong hết sức thuyết phục~\cite{ji2023hallucination, huang2023hallucination}, hoặc đưa ra nội dung không bám sát chuẩn kiến thức và kỹ năng của chương trình giáo dục phổ thông quốc gia. Mặt khác, tri thức nội tại của mô hình bị đóng băng tại thời điểm hoàn tất quá trình huấn luyện; các hệ thống trợ lý hội thoại thông thường hiện nay thiếu hẳn cơ chế truy xuất tường minh từ các nguồn tài liệu sư phạm chính thống như sách giáo khoa chuẩn mực hay học liệu đã qua thẩm định chuyên môn~\cite{lewis2020retrieval, gao2023ragsurvey}. Sự thiếu vắng khả năng đối chiếu và kiểm chứng nguồn gốc khiến độ tin cậy của câu trả lời bị suy giảm nghiêm trọng, gây trở ngại lớn cho việc ứng dụng công nghệ này trong môi trường giáo dục chính quy.

Để khắc phục triệt để các hạn chế trên, hướng tiếp cận sinh tăng cường truy xuất RAG~\cite{lewis2020retrieval} đã khẳng định vị thế là một giải pháp tiên tiến, kết hợp hài hòa năng lực lập luận ngôn ngữ linh hoạt của mô hình ngôn ngữ lớn với khả năng truy xuất tri thức chuẩn xác từ kho tài liệu chuyên biệt và đáng tin cậy~\cite{gao2023ragsurvey, yan2024crag, jeong2024adaptiverag}. Thay vì phụ thuộc hoàn toàn vào các trọng số tham số tĩnh, mô hình RAG chủ động tìm kiếm các đoạn thông tin liên quan nhất từ cơ sở tri thức xác thực để làm ngữ cảnh đầu vào cho mô hình ngôn ngữ tổng hợp câu trả lời hoàn chỉnh~\cite{lewis2020retrieval, guu2020realm}. Phương pháp này nâng cao vượt bậc tính chuẩn xác của câu trả lời, đảm bảo tính giải trình khoa học thông qua việc dẫn xuất tường minh nguồn gốc của từng luận điểm, đồng thời cho phép cập nhật tri thức linh hoạt mà không phát sinh chi phí tái huấn luyện mô hình tốn kém~\cite{gao2023ragsurvey}.

Đặc biệt, đặc thù nổi bật của sách giáo khoa môn Khoa học tự nhiên cấp Trung học cơ sở là nội dung lý thuyết luôn gắn liền mật thiết với hệ thống \textbf{hình ảnh trực quan, sơ đồ giải phẫu, biểu đồ thí nghiệm và phương trình khoa học}: từ cấu trúc tế bào sinh học, sơ đồ mạch điện, đồ thị chuyển động cho đến các phản ứng hoá học. Một hệ thống RAG truyền thống thuần văn bản sẽ bỏ sót toàn bộ kho tài nguyên trực quan quý giá này. Do đó, việc nghiên cứu phát triển một kiến trúc RAG \textbf{đa phương thức}, có năng lực thấu hiểu và truy xuất đồng thời cả dữ liệu văn bản lẫn hình ảnh dựa trên các mô hình liên kết thị giác -- ngôn ngữ hiện đại như CLIP~\cite{radford2021clip}, là một yêu cầu khoa học cấp thiết.

Bên cạnh bài toán mở rộng không gian tri thức đa phương thức, \textbf{thách thức trọng tâm trong nghiên cứu RAG nằm ở chất lượng của pha truy xuất}. Trong thực tế triển khai, một tỷ lệ đáng kể các đoạn tài liệu chính xác thường bị loại bỏ ở khâu chấm điểm và chọn lọc ngưỡng $top\text{-}k$, thay vì do không gian nhúng không thể tìm thấy chúng. Hệ quả là độ phủ ở quy mô vận hành thực tế thường thấp hơn nhiều so với năng lực trần lý thuyết của riêng bước tìm kiếm ngữ nghĩa nếu mở rộng số lượng đoạn ứng viên. Chính khoảng cách này phản ánh giới hạn có thể tối ưu hoá bằng phương pháp luận kiến trúc chứ không chỉ đơn thuần là tăng quy mô tham số mô hình. Cần nhấn mạnh rằng, các kết quả thực nghiệm định lượng tại Chương~\ref{chapter:thuc_nghiem_danh_gia} sẽ phân tích chi tiết mức độ cải thiện cũng như những giới hạn thực tế của khoảng cách này (xem Mục~\ref{sec:ablation}). Vì vậy, đề tài tập trung vào hai hướng nghiên cứu then chốt: (i) nâng cao độ tin cậy của luồng truy xuất đa phương thức thông qua cơ chế định tuyến ý định, truy xuất lai và bước tái xếp hạng chuyên sâu, và (ii) \textbf{kiểm chứng thực nghiệm đối chiếu khách quan}: so sánh truy xuất lai với mô hình BM25 truyền thống~\cite{robertson2009bm25}, và đối chiếu RAG đa phương thức với RAG thuần văn bản.

Xuất phát từ những yêu cầu thực tiễn và thách thức nghiên cứu nêu trên, đề tài \textbf{"Nghiên cứu phương pháp RAG đa phương thức cho trợ lý ảo thông minh hỗ trợ dạy và học môn Khoa học tự nhiên cấp Trung học cơ sở"} được thực hiện nhằm xây dựng một giải pháp trợ lý học tập thông minh, đảm bảo cung cấp học liệu số chuẩn xác, minh bạch và có khả năng truy nguyên chính xác về từng trang sách giáo khoa phục vụ học sinh và giáo viên.

\section{Mục tiêu của đề tài}
Mục tiêu bao trùm của đề tài là nghiên cứu, thiết kế, hiện thực hóa và đánh giá định lượng phương pháp RAG đa phương thức ứng dụng trong trợ lý ảo hỗ trợ dạy và học môn Khoa học tự nhiên cấp Trung học cơ sở. Các mục tiêu nghiên cứu cụ thể bao gồm:

\begin{itemize}
    \item \textbf{Mục tiêu 1 --- Xây dựng cơ chế truy vấn đa định dạng:} Xây dựng mô hình cho phép hệ thống tiếp nhận, thấu hiểu và truy xuất đồng thời trên nhiều định dạng nội dung đặc thù của môn Khoa học tự nhiên --- văn bản lý thuyết, \textbf{công thức khoa học}, hình ảnh minh hoạ, sơ đồ và biểu đồ thí nghiệm --- khắc phục triệt để hạn chế của các hệ thống thuần văn bản.
    \item \textbf{Mục tiêu 2 --- Thiết lập cơ sở tri thức vector toàn diện:} Thực hiện số hoá, tiền xử lý và lập chỉ mục \textbf{toàn bộ} chương trình môn Khoa học tự nhiên từ lớp 6 đến lớp 9 của cả ba bộ sách giáo khoa hiện hành gồm Kết nối tri thức với cuộc sống, Chân trời sáng tạo và Cánh Diều~\cite{sgk_kntt6, sgk_kntt7, sgk_kntt8, sgk_kntt9, sgk_ctst6, sgk_ctst7, sgk_ctst8, sgk_ctst9, sgk_cd6, sgk_cd7, sgk_cd8, sgk_cd9}; xử lý chuyên sâu cả văn bản lý thuyết, ký hiệu công thức và hình ảnh minh họa.
    \item \textbf{Mục tiêu 3 --- Tối ưu hóa pha truy xuất tri thức:} Hoàn thiện cơ chế định tuyến theo ý định truy vấn, phát triển mô hình truy xuất lai kết hợp giữa tìm kiếm từ khoá BM25 và tìm kiếm ngữ nghĩa vector, cùng bước tái xếp hạng bằng mô hình tương tác chéo, nhằm \textbf{thu hẹp khoảng cách} giữa độ phủ thực tế và trần độ phủ lý thuyết của hệ thống.
    \item \textbf{Mục tiêu 4 --- Thiết lập khung đánh giá thực nghiệm đối chiếu:} Xây dựng quy trình đánh giá định lượng chặt chẽ và thực hiện \textbf{thực nghiệm so sánh đối đầu} trên cùng một bộ dữ liệu kiểm thử: (i) truy xuất lai so với mô hình BM25 truyền thống, (ii) RAG đa phương thức so với RAG thuần văn bản; đồng thời tiến hành thực nghiệm bóc tách thành phần nhằm phân tích vai trò của bước tái xếp hạng và cổng lọc theo độ liên quan.
    \item \textbf{Mục tiêu 5 --- Hiện thực hóa ứng dụng trợ lý học tập trực quan:} Phát triển hoàn chỉnh ứng dụng web tương tác cho phép giáo viên và học sinh đặt câu hỏi và nhận câu trả lời kèm hình ảnh trích xuất, \textbf{công thức hiển thị chuẩn ký hiệu khoa học} và thông tin trích dẫn minh bạch đến từng trang sách giáo khoa nguồn.
    \item \textbf{Cơ sở phương pháp luận cho các mục tiêu --- Bộ dữ liệu kiểm thử chuẩn đối sánh:} Xây dựng bộ dữ liệu kiểm thử được gắn nhãn nguồn xác định theo tên sách và số trang làm chuẩn đối sánh khách quan, bao quát cả ba phân môn, bốn khối lớp và ba bộ sách giáo khoa. Pha truy xuất được lượng hoá bằng các chỉ số Precision@$k$, Recall@$k$ và MRR đối chiếu trực tiếp nhãn nguồn; pha sinh ngôn ngữ được đánh giá qua phương pháp LLM-as-a-judge~\cite{zheng2023llmjudge} với mô hình ngôn ngữ độc lập chấm điểm theo ba tiêu chí cốt lõi: Tính đúng, Độ trung thực và Độ liên quan, dựa trên các khung tiêu chuẩn đo lường RAG chuẩn hóa~\cite{es2024ragas, saadfalcon2023ares}.
\end{itemize}

\section{Đối tượng và phạm vi nghiên cứu}
\subsection{Đối tượng nghiên cứu}
\begin{itemize}
    \item Các phương pháp xử lý ngôn ngữ tự nhiên và trí tuệ nhân tạo tạo sinh.
    \item Các mô hình ngôn ngữ lớn mã nguồn mở có năng lực tuân thủ chỉ thị cao.
    \item Kiến trúc hệ thống sinh tăng cường truy xuất đa phương thức RAG xử lý đồng thời dữ liệu văn bản và hình ảnh.
    \item Cơ sở dữ liệu vector, các mô hình nhúng ngữ nghĩa và mô hình tái xếp hạng tương tác chéo.
\end{itemize}

\subsection{Phạm vi nghiên cứu}
\begin{itemize}
    \item \textbf{Về phạm vi dữ liệu:} Nghiên cứu bao quát \textbf{toàn bộ} chương trình môn Khoa học tự nhiên cấp Trung học cơ sở từ lớp 6 đến lớp 9 theo Chương trình Giáo dục phổ thông 2018, tích hợp cả ba phân môn \textbf{Vật lí, Hoá học và Sinh học}. Cơ sở tri thức được xây dựng từ 12 cuốn sách giáo khoa thuộc ba bộ sách hiện hành: Kết nối tri thức với cuộc sống~\cite{sgk_kntt6, sgk_kntt7, sgk_kntt8, sgk_kntt9}, Chân trời sáng tạo~\cite{sgk_ctst6, sgk_ctst7, sgk_ctst8, sgk_ctst9} và Cánh Diều~\cite{sgk_cd6, sgk_cd7, sgk_cd8, sgk_cd9}. Do đó, các nội dung thuộc phân môn Vật lí và Hoá học là thành phần cốt lõi trong phạm vi nghiên cứu của đề tài.
    \item \textbf{Về mặt kỹ thuật:} Hệ thống được thiết kế theo kiến trúc máy khách -- máy chủ hiện đại. Phía máy chủ chịu trách nhiệm quản lý quy trình RAG, điều phối mô hình nhúng và mô hình ngôn ngữ lớn sử dụng ngôn ngữ Python, thư viện LangChain và cơ sở dữ liệu vector ChromaDB. Phía máy khách cung cấp giao diện tương tác trực quan cho người dùng phát triển trên nền tảng Next.js.
    \item \textbf{Về đối tượng phục vụ:} Học sinh và giáo viên cấp Trung học cơ sở có nhu cầu tra cứu học liệu, ôn tập kiến thức và giải đáp thắc mắc liên quan đến chương trình môn Khoa học tự nhiên.
    \item \textbf{Giới hạn phạm vi nghiên cứu:} Đề tài không hướng đến việc tinh chỉnh toàn bộ trọng số của các mô hình ngôn ngữ lớn, mà tập trung chủ đạo vào việc hoàn thiện cấu trúc và nâng cao độ chính xác của quy trình truy xuất tri thức cùng khả năng tích hợp đa phương thức.
\end{itemize}

\section{Cấu trúc của báo cáo đồ án}
Báo cáo đồ án tốt nghiệp được tổ chức thành 5 chương với nội dung tóm tắt như sau:

\begin{itemize}
    \item \textbf{Chương 1: Tổng quan đề tài.} Trình bày bối cảnh thực tiễn, tính cấp thiết, mục tiêu nghiên cứu, đối tượng, phạm vi và bố cục tổng thể của báo cáo.
    \item \textbf{Chương 2: Cơ sở lý thuyết.} Hệ thống hóa các kiến thức nền tảng phục vụ nghiên cứu, bao gồm các nguyên lý xử lý ngôn ngữ tự nhiên, mô hình ngôn ngữ lớn, kiến trúc RAG, hệ quản trị cơ sở dữ liệu vector, kỹ thuật biểu diễn đa phương thức và các phương pháp truy xuất kết hợp tái xếp hạng.
    \item \textbf{Chương 3: Phương pháp thực hiện và kiến trúc hệ thống.} Mô tả chi tiết kiến trúc tổng thể của hệ thống trợ lý ảo và các công nghệ nền tảng. Trình bày quy trình trích xuất và tiền xử lý dữ liệu cho văn bản, công thức và hình ảnh; cách thức thiết kế luồng RAG đa phương thức và giải pháp giao diện người dùng.
    \item \textbf{Chương 4: Thực nghiệm, đánh giá và thảo luận.} Trình bày môi trường thực nghiệm, quy trình xây dựng bộ dữ liệu kiểm thử gắn nhãn nguồn khách quan, hệ thống chỉ số đánh giá hai pha truy xuất và sinh ngôn ngữ; phân tích chi tiết kết quả định lượng qua các bảng đối chiếu thực nghiệm, thảo luận chuyên sâu về các phát hiện khoa học và nhận diện các giới hạn kỹ thuật của hệ thống.
    \item \textbf{Chương 5: Kết luận và hướng phát triển.} Tổng kết các kết quả nghiên cứu và mức độ hoàn thành so với các mục tiêu đề ra, đồng thời đề xuất lộ trình và giải pháp phát triển hệ thống trong tương lai.
\end{itemize}

\chapter{CƠ SỞ LÝ THUYẾT}
\label{chapter:co_so_ly_thuyet}

Chương này hệ thống hóa các cơ sở lý thuyết và nguyên lý khoa học nền tảng làm tiền đề cho việc thiết kế, xây dựng và đánh giá hệ thống trợ lý học tập thông minh. Nội dung trọng tâm bao gồm: tổng quan về xử lý ngôn ngữ tự nhiên và trí tuệ nhân tạo tạo sinh, kiến trúc và cơ chế hoạt động của các mô hình ngôn ngữ lớn, phương pháp sinh tăng cường truy xuất RAG, hệ quản trị cơ sở dữ liệu vector, các kỹ thuật biểu diễn đa phương thức, cùng các giải pháp truy xuất lai và kỹ thuật tái xếp hạng tương tác chéo.

\section{Tổng quan về Xử lý ngôn ngữ tự nhiên và Trí tuệ nhân tạo tạo sinh}
Xử lý ngôn ngữ tự nhiên là một phân ngành nền tảng của trí tuệ nhân tạo, hướng tới mục tiêu phát triển các phương pháp và mô hình tính toán cho phép máy tính có khả năng tiếp nhận, giải mã cú pháp, thấu hiểu ngữ nghĩa và tạo lập ngôn ngữ con người một cách tự nhiên. Lịch sử phát triển của lĩnh vực này ghi nhận sự dịch chuyển mang tính bước ngoặt: từ các hệ thống dựa trên tập luật hình thức thủ công và các mô hình thống kê tần suất sang các kiến trúc học sâu hiện đại~\cite{bengio2003neural, lecun2015deeplearning}. Bước đột phá khởi đầu cho giai đoạn này là sự xuất hiện của các phương pháp biểu diễn ngữ nghĩa phân tán trong không gian vector liên tục, tiêu biểu như mô hình word2vec~\cite{mikolov2013word2vec}, từ đó mở đường cho các kiến trúc mạng nơ-ron hồi quy và mô hình biến đổi chuỗi sang chuỗi tích hợp cơ chế chú ý phục vụ dịch máy tự động~\cite{sutskever2014seq2seq, bahdanau2015attention}.

Trên cơ sở đó, sự bùng nổ của trí tuệ nhân tạo tạo sinh đã thiết lập một bước chuyển biến sâu sắc về năng lực ứng dụng thực tiễn. Khác với các mô hình phân biệt truyền thống vốn chỉ tập trung ước lượng xác suất có điều kiện nhằm phân loại hoặc nhận dạng các mẫu dữ liệu có sẵn, các mô hình tạo sinh học phân phối xác suất đồng thời trên tập dữ liệu huấn luyện, cho phép tổng hợp và khởi tạo những đối tượng dữ liệu hoàn toàn mới mang tính nhất quán cao về mặt ngữ nghĩa, bao gồm văn bản, hình ảnh, âm thanh và mã nguồn phần mềm~\cite{bommasani2021foundation}. Trong không gian dữ liệu văn bản, nhân tố đóng vai trò hạt nhân thúc đẩy cuộc cách mạng này chính là các mô hình ngôn ngữ lớn~\cite{zhao2023llmsurvey}.

\section{Mô hình ngôn ngữ lớn}
\subsection{Kiến trúc Transformer và cơ chế tự chú ý}
Được công bố bởi Vaswani và các cộng sự vào năm 2017~\cite{vaswani2017attention}, kiến trúc Transformer đã nhanh chóng thay thế các mô hình tuần tự truyền thống để trở thành cấu trúc nền tảng cho hầu hết các mô hình ngôn ngữ lớn hiện đại. Trọng tâm của Transformer là cơ chế tự chú ý đa điểm, cho phép mô hình tính toán đồng thời trọng số tương tác giữa từng từ tố với toàn bộ các từ tố còn lại trong ngữ cảnh, hoàn toàn không bị giới hạn bởi khoảng cách vị trí giữa chúng.

Cơ chế này khắc phục triệt để hiện tượng suy giảm gradient và hạn chế mất mát thông tin ngữ cảnh xa vốn xuất hiện phổ biến ở các mạng hồi quy như RNN hay LSTM. Đồng thời, kiến trúc Transformer giải phóng năng lực tính toán song song quy mô lớn trên các bộ xử lý đồ họa, tạo điều kiện tiên quyết cho việc mở rộng quy mô dữ liệu và tham số huấn luyện lên tới hàng trăm tỷ trọng số. Từ cấu trúc Transformer, hai nhánh mô hình tiền huấn luyện chủ đạo đã phát triển mạnh mẽ: nhánh bộ mã hóa hai chiều tiêu biểu như BERT~\cite{devlin2019bert} tối ưu cho các bài toán phân loại và thấu hiểu ngữ nghĩa; và nhánh bộ giải mã tự hồi quy như dòng GPT~\cite{brown2020gpt3} chuyên biệt cho tác vụ tạo lập văn bản nối tiếp. Các nghiên cứu định lượng đã khẳng định rằng khi mở rộng quy mô tham số, thể tích dữ liệu và năng lực tính toán, chất lượng mô hình tuân theo các quy luật tỷ lệ chặt chẽ~\cite{kaplan2020scaling}, làm bộc lộ những năng lực suy luận phức tạp mà các mô hình quy mô nhỏ không thể đạt được~\cite{wei2022emergent}.

\subsection{Họ mô hình ngôn ngữ lớn Qwen}
Trong số các dòng mô hình ngôn ngữ lớn mã nguồn mở hiện nay, họ mô hình Qwen, tiêu biểu là thế hệ Qwen2.5 do Alibaba Cloud phát triển~\cite{qwen25_2024}, được đánh giá là một trong những giải pháp dẫn đầu về hiệu năng tổng thể. Sự phát triển mạnh mẽ của các dòng mô hình mã nguồn mở như LLaMA~\cite{touvron2023llama} hay Qwen mang lại cho cộng đồng nghiên cứu khả năng tự chủ công nghệ cao, cho phép triển khai hoàn toàn cục bộ, bảo mật dữ liệu và tùy biến chuyên sâu mà không bị lệ thuộc vào các giao diện lập trình thương mại đóng kín.

Trong khuôn khổ đồ án, phiên bản \textbf{Qwen2.5-3B-Instruct} được lựa chọn làm mô hình sinh ngôn ngữ nòng cốt. Mô hình được tinh chỉnh chuyên sâu theo chỉ thị kết hợp phương pháp học tăng cường từ phản hồi của con người~\cite{ouyang2022instructgpt}, đem lại năng lực tuân thủ chỉ thị và lập luận vượt trội trong phân khúc mô hình nhỏ gọn với quy mô khoảng 3 tỷ tham số. Kích thước tối ưu này cho phép hệ thống đáp ứng tốt các bài toán thực tế trên cơ sở hạ tầng có tài nguyên tính toán giới hạn mà vẫn bảo đảm tốc độ suy luận nhanh và chất lượng văn bản sinh ra mạch lạc. Toàn bộ quy trình thực nghiệm và kiểm thử của đồ án được triển khai thống nhất trên máy trạm chỉ sử dụng \textbf{CPU} (chi tiết tại Bảng~\ref{tab:env}). Đặc biệt, mô hình Qwen2.5 thể hiện năng lực đa ngữ xuất sắc, thấu hiểu và diễn đạt tiếng Việt trôi chảy, chuẩn xác, đáp ứng tốt yêu cầu giải thích kiến thức khoa học mang tính sư phạm trong môi trường giáo dục Việt Nam~\cite{nguyen2020phobert}.

\section{Sinh tăng cường truy xuất RAG}
\subsection{Những hạn chế cố hữu của mô hình ngôn ngữ lớn độc lập}
Mặc dù thể hiện khả năng biểu đạt ngôn ngữ rất tự nhiên, các mô hình ngôn ngữ lớn khi triển khai độc lập đều bộc lộ ba điểm nghẽn kỹ thuật căn bản~\cite{gao2023ragsurvey}:
\begin{enumerate}
    \item \textbf{Ảo giác tri thức:} Mô hình có thể tạo lập các khẳng định sai lệch hoặc suy đoán thiếu cơ sở thực tế với văn phong hết sức tự tin, thiếu sự đối sánh với dữ liệu xác thực~\cite{ji2023hallucination, huang2023hallucination}.
    \item \textbf{Tính chất tri thức tĩnh:} Toàn bộ tri thức của mô hình bị đóng băng tại thời điểm hoàn tất tập dữ liệu huấn luyện, không thể tự cập nhật thông tin mới trừ khi phải trải qua quá trình tái huấn luyện đòi hỏi chi phí tính toán vô cùng đắt đỏ~\cite{lewis2020retrieval}.
    \item \textbf{Sự thiếu hụt tri thức chuyên biệt:} Mô hình không có sẵn quyền truy cập vào các kho học liệu nội bộ, tài liệu đặc thù của tổ chức hay nội dung chi tiết của các bộ sách giáo khoa chuẩn mực quốc gia~\cite{guu2020realm}.
\end{enumerate}

\subsection{Nguyên lý hoạt động của kiến trúc RAG}
\label{sec:rag_nguyenly}
Kiến trúc sinh tăng cường truy xuất RAG~\cite{lewis2020retrieval} ra đời nhằm giải quyết triệt để các hạn chế trên thông qua sự kết hợp cộng sinh giữa hai thành phần: hệ thống truy xuất thông tin và mô hình tạo sinh ngôn ngữ. Thay vì xem mô hình ngôn ngữ như một kho lưu trữ tri thức tham số khép kín, mô hình RAG tách biệt thành phần ghi nhớ tri thức sang một cơ sở dữ liệu bên ngoài, xem mô hình ngôn ngữ lớn như một bộ xử lý suy luận hoạt động trên ngữ cảnh được cung cấp tức thời~\cite{guu2020realm, izacard2021fid, gao2023ragsurvey, asai2023selfrag, yan2024crag}.

\begin{figure}[H]
    \centering
    \includegraphics[width=\textwidth]{src/images/chapter2/rag_concept.png}
    \caption{Sơ đồ nguyên lý hoạt động cơ bản của hệ thống RAG (Nguồn: Tổng hợp từ~\cite{lewis2020retrieval, gao2023ragsurvey})}
    \label{fig:rag_concept}
\end{figure}

Quy trình vận hành chuẩn của một hệ thống RAG bao gồm hai pha chính (Hình~\ref{fig:rag_concept}):
\begin{itemize}
    \item \textbf{Pha Lập chỉ mục:} Dữ liệu tài liệu thô từ các tệp PDF và ảnh chụp trang sách được trích xuất và phân chia thành các đoạn văn bản độc lập. Tiếp đó, mô hình nhúng chuyển đổi các đoạn này thành các vector số học đa chiều để lưu trữ và lập chỉ mục trong cơ sở dữ liệu vector.
    \item \textbf{Pha Truy vấn và Tạo sinh:} Khi người dùng gửi câu hỏi, hệ thống tiến hành nhúng câu hỏi thành vector truy vấn tương ứng. Cơ chế tìm kiếm độ tương đồng trong cơ sở dữ liệu vector sẽ xác định các đoạn văn bản có độ tương đồng ngữ nghĩa cao nhất để trích xuất thành tập ngữ cảnh liên quan. Phương pháp này, hay còn được gọi là truy xuất ngữ nghĩa theo vector đặc~\cite{karpukhin2020dpr}, cho phép tìm kiếm chính xác nội dung đồng nghĩa ngay cả khi không có sự trùng khớp hoàn toàn về mặt từ khóa, khắc phục nhược điểm của các thuật toán so khớp từ khóa thuần túy như BM25~\cite{robertson2009bm25}. Cuối cùng, lời nhắc có cấu trúc kết hợp giữa ngữ cảnh tài liệu và câu hỏi gốc sẽ được đưa vào mô hình ngôn ngữ để tạo lập câu trả lời hoàn chỉnh có căn cứ học liệu rõ ràng.
\end{itemize}

\section{Mô hình nhúng và Cơ sở dữ liệu Vector}
\subsection{Kỹ thuật biểu diễn nhúng ngữ nghĩa}
Kỹ thuật nhúng là quá trình ánh xạ các thực thể ngôn ngữ như từ, câu, đoạn văn bản hoặc phi ngôn ngữ như hình ảnh vào một không gian vector thực liên tục đa chiều $\mathbb{R}^d$~\cite{mikolov2013word2vec}. Đặc tính quan trọng nhất của không gian nhúng là bảo toàn quan hệ ngữ nghĩa: các thực thể có nội dung hoặc ý nghĩa tương đồng sẽ nằm gần nhau trong không gian vector đa chiều, được định lượng thông qua các phép đo như độ tương đồng cosine. Đối với cấp độ câu và đoạn văn bản, các kiến trúc dựa trên Sentence-BERT~\cite{reimers2019sbert} cho phép trích xuất các vector biểu diễn ngữ nghĩa cô đọng và có khả năng so khớp trực tiếp với hiệu quả tính toán vượt trội.

Trong hệ thống nghiên cứu, mô hình \texttt{BAAI/bge-m3}~\cite{chen2024bgem3} được sử dụng để sinh các vector biểu diễn ngữ nghĩa với kích thước $1024$ chiều. Đây là mô hình đa ngữ tiên tiến được thiết kế hỗ trợ đồng thời ba cơ chế biểu diễn gồm vector đặc, vector thưa và đa vector biểu diễn, đồng thời có khả năng tiếp nhận độ dài ngữ cảnh lớn lên tới $8192$ token, rất phù hợp với đặc thù của các đoạn văn bản sách giáo khoa thường đi kèm công thức, ký hiệu khoa học và thuật ngữ định nghĩa.

So với các kiến trúc nhúng nông quy mô $384$ chiều như \texttt{MiniLM}~\cite{wang2020minilm}, mô hình \texttt{BAAI/bge-m3} thể hiện năng lực biểu diễn ngữ nghĩa vượt trội trên các bảng xếp hạng chuẩn MTEB~\cite{muennighoff2022mteb} và BEIR~\cite{thakur2021beir}, đặc biệt đối với các tác vụ truy xuất học liệu đa ngữ phức tạp. Về mặt nguyên lý toán học của không gian vector, các mô hình nhúng có số chiều khác nhau sẽ định hình các không gian mêtric hoàn toàn tách biệt, không thể chiếu trực tiếp hay tính toán độ tương đồng qua lại. Do đó, việc thiết lập một không gian vector $1024$ chiều đồng nhất xuyên suốt toàn bộ kho dữ liệu là yêu cầu tiên quyết nhằm đảm bảo tính chuẩn xác cho các hàm khoảng cách hình học.

\subsection{Hệ quản trị cơ sở dữ liệu vector ChromaDB}
Cơ sở dữ liệu vector là hệ quản trị dữ liệu chuyên dụng được tối ưu hóa cho bài toán lưu trữ và tìm kiếm vector lân cận gần đúng trên quy mô dữ liệu lớn. Khác với các cơ sở dữ liệu quan hệ truyền thống chỉ hỗ trợ so khớp chính xác, cơ sở dữ liệu vector cho phép thực hiện truy vấn tương đồng tốc độ cao dựa trên các chỉ mục đồ thị lân cận phân tầng HNSW~\cite{malkov2018hnsw} hoặc các thư viện tính toán tối ưu như FAISS~\cite{johnson2019faiss}.

Trong đề tài này, ChromaDB~\cite{chromadb2023} được lựa chọn làm cơ sở dữ liệu vector trung tâm. ChromaDB là một giải pháp mã nguồn mở linh hoạt, hỗ trợ vận hành cục bộ ở chế độ nhúng với hiệu năng cao mà không yêu cầu thiết lập máy chủ dịch vụ phức tạp, đồng thời tích hợp liền mạch với hệ sinh thái LangChain~\cite{langchain2022}, tạo điều kiện thuận lợi cho việc tổ chức, cập nhật và truy xuất các phân mảnh tri thức từ sách giáo khoa.

\section{Kỹ thuật xử lý Đa phương thức trong kiến trúc RAG}
Nội dung sách giáo khoa môn Khoa học tự nhiên chứa mật độ cao các đối tượng trực quan: sơ đồ cấu tạo cơ quan sinh học, chu trình chuyển hóa, sơ đồ mạch điện, đồ thị dao động, thiết bị thí nghiệm... Một hệ thống RAG thuần văn bản sẽ bỏ lỡ hoàn toàn các tri thức trực quan này. Do đó, việc mở rộng sang kiến trúc đa phương thức là một đòi hỏi tất yếu, dựa trên nền tảng của các mô hình Vision Transformer~\cite{dosovitskiy2021vit} và các kiến trúc liên kết ngôn ngữ -- thị giác hiện đại~\cite{radford2021clip, li2022blip, liu2024llava, li2023blip2, alayrac2022flamingo}, kế thừa các thành tựu từ bài toán hỏi đáp trực quan trên hình ảnh~\cite{antol2015vqa} và xử lý tài liệu đa phương thức~\cite{cui2021documentai, huang2022layoutlmv3}.

\subsection{Các mô hình trích xuất và biểu diễn thị giác}
Hệ thống tích hợp có chọn lọc các mô hình thị giác máy tính chuyên dụng:
\begin{itemize}
    \item \textbf{Mô hình OwlViT}\footnote{\texttt{google/owlvit-base-patch32}}~\cite{minderer2022owlvit}: Mô hình phát hiện vật thể không cần huấn luyện trước dựa trên điều kiện văn bản, được ứng dụng nhằm hỗ trợ xác định toạ độ hộp bao của các biểu đồ và hình ảnh trên trang sách trong các trường hợp không có nhãn văn bản cố định.
    \item \textbf{Mô hình CLIP}\footnote{\texttt{openai/clip-vit-base-patch16}}~\cite{radford2021clip}: Mô hình liên kết ngôn ngữ -- thị giác được tiền huấn luyện trên hàng trăm triệu cặp ảnh--văn bản. CLIP mã hoá hình ảnh thành vector $512$ chiều chung một không gian biểu diễn với câu hỏi văn bản, cho phép thực hiện tìm kiếm chéo giữa các phương thức thông qua độ tương đồng cosine.
    \item \textbf{Mô hình Vintern-1B}\footnote{\texttt{5CD-AI/Vintern-1B-v2}}: Mô hình thị giác--ngôn ngữ hỗ trợ tiếng Việt có khả năng tự động tạo chuỗi miêu tả văn bản từ hình ảnh cắt ra. Trong đề tài, mô hình này đã được đưa vào thử nghiệm thực tế; tuy nhiên, do vấn đề phát sinh ảo giác và chi phí tính toán cao trên CPU, hệ thống đã lựa chọn giải pháp trích xuất siêu dữ liệu trực tiếp từ văn bản gốc thay vì dựa vào chú thích tự sinh (trình bày chi tiết tại Chương~\ref{chapter:phuong_phap_thuc_hien}).
\end{itemize}

Thông qua sự kết hợp của các công nghệ trên, kiến trúc RAG đa phương thức của hệ thống đảm bảo khả năng đồng bộ hóa luồng tri thức giữa văn bản lý thuyết và dữ liệu hình ảnh, tạo tiền đề cung cấp ngữ cảnh học liệu toàn diện nhất cho mô hình sinh câu trả lời.

\section{Truy xuất lai và Kỹ thuật tái xếp hạng}
\label{sec:cosolythuyet_hybrid}
Truy xuất ngữ nghĩa dựa trên vector đặc tuy có năng lực nắm bắt ngữ cảnh tốt nhưng không phải là phương pháp tối ưu tuyệt đối trong mọi tình huống, đặc biệt là khi phải xử lý các truy vấn chứa ký hiệu chuyên ngành hoặc thuật ngữ cụ thể. Nhằm đạt được hiệu năng truy xuất tối ưu, hệ thống kết hợp ba kỹ thuật nền tảng sau:

\subsection{Truy xuất từ khóa theo vector thưa: Thuật toán BM25}
Okapi BM25~\cite{robertson2009bm25} là thuật toán xếp hạng tài liệu dựa trên mô hình xác suất thống kê tần suất từ khóa, là chuẩn mực trong lĩnh vực truy xuất thông tin cổ điển. Với một câu truy vấn $Q$ bao gồm các từ tố $\{q_1, q_2, \dots, q_n\}$ và một tài liệu $D$, hàm điểm số BM25 được xác định theo công thức:
\begin{equation}
\mathrm{BM25}(D, Q) = \sum_{i=1}^{n} \mathrm{IDF}(q_i) \cdot
\frac{f(q_i, D)\,(k_1 + 1)}{f(q_i, D) + k_1\left(1 - b + b\,\dfrac{|D|}{\mathrm{avgdl}}\right)}
\end{equation}
trong đó $f(q_i, D)$ biểu thị tần suất xuất hiện của từ $q_i$ trong tài liệu $D$; $|D|$ là độ dài tài liệu; $\mathrm{avgdl}$ là độ dài trung bình của toàn bộ các tài liệu trong tập dữ liệu; $\mathrm{IDF}(q_i)$ là nghịch số tần suất tài liệu của từ $q_i$; $k_1$ và $b$ là hai siêu tham số thực nghiệm kiểm soát mức độ bão hòa tần suất và mức độ chuẩn hóa theo độ dài tài liệu.

Điểm mạnh cốt lõi của BM25 là khả năng so khớp chính xác từng ký tự. Đối với môn Khoa học tự nhiên vốn chứa mật độ cao các công thức hóa học như \ce{H2SO4} hay \ce{Fe2O3}, ký hiệu vật lí hoặc tên riêng của các định luật, đây là một lợi thế quyết định: câu hỏi của người học đòi hỏi tài liệu trích xuất phải chứa chính xác chuỗi ký hiệu đó mà không bị pha loãng bởi phép xấp xỉ ngữ nghĩa. Ngược lại, hạn chế của BM25 là không thể nhận diện được các từ đồng nghĩa hoặc các cách diễn đạt khác nhau cho cùng một khái niệm khoa học.

\subsection{Cơ chế truy xuất lai và hợp nhất thứ hạng}
Do kênh truy xuất thưa BM25 và kênh truy xuất vector đặc sở hữu các đặc tính bổ trợ cho nhau, hướng tiếp cận hiệu quả là kết hợp song song cả hai kênh rồi hợp nhất các danh sách kết quả xếp hạng. Một phương pháp hợp nhất được ứng dụng rộng rãi là phương pháp hợp nhất theo thứ hạng nghịch đảo RRF~\cite{cormack2009reciprocal}. Trong phương pháp RRF, điểm số hợp nhất của mỗi tài liệu $d$ được tính toán dựa trên vị trí thứ hạng của nó trong các bảng xếp hạng thành phần:
\begin{equation}
\mathrm{RRF}(d) = \sum_{r \in R} \frac{1}{k + \mathrm{rank}_r(d)}
\end{equation}
với $R$ là tập hợp các bảng xếp hạng từ các kênh truy xuất thành phần, $\mathrm{rank}_r(d)$ là vị trí thứ hạng của tài liệu $d$ trong kênh $r$, và $k$ là hằng số làm mượt thực nghiệm với giá trị thường chọn là $k = 60$. Ưu thế vượt trội của RRF là chỉ dựa vào thứ hạng tương đối thay vì điểm số tuyệt đối, do đó loại bỏ hoàn toàn yêu cầu phải chuẩn hoá thang đo điểm số vốn rất khác biệt giữa khoảng cách cosine trong không gian vector đặc và điểm số xác suất mở của BM25.

Bảng~\ref{tab:so_sanh_fusion} so sánh các phương pháp hợp nhất kết quả truy xuất lai phổ biến trong các hệ thống truy xuất thông tin hiện đại.

\begin{table}[H]
\centering
\caption{So sánh các phương pháp hợp nhất kết quả trong truy xuất lai (Nguồn: Tác giả tổng hợp)}
\label{tab:so_sanh_fusion}
\small
\begin{tabularx}{\textwidth}{>{\raggedright\arraybackslash\bfseries}p{3.0cm} >{\raggedright\arraybackslash}X >{\raggedright\arraybackslash}X >{\raggedright\arraybackslash}X}
\toprule
\textbf{Phương pháp} & \textbf{Nguyên lý tính toán} & \textbf{Ưu điểm} & \textbf{Hạn chế} \\
\midrule
Hợp nhất thứ hạng nghịch đảo RRF & Dựa trên thứ bậc tương đối $\sum \frac{1}{k + \mathrm{rank}}$ & Không phụ thuộc phân phối điểm số; không cần tham số học & Bỏ qua khoảng cách điểm tuyệt đối giữa các vị trí \\
Chuẩn hóa Min-Max kết hợp trọng số & Đưa điểm hai kênh về $[0, 1]$ rồi cộng trọng số & Giữ lại độ lớn điểm tương đồng ban đầu & Nhạy cảm với giá trị ngoại lai trong từng lô \\
Kết hợp lồi & $S = \alpha S_{\text{dense}} + (1-\alpha) S_{\text{sparse}}$ & Dễ tinh chỉnh tham số $\alpha$ theo từng miền dữ liệu & Đòi hỏi chuẩn hóa phức tạp; hiệu năng phụ thuộc tập dữ liệu \\
Tương tác muộn ColBERT~\cite{khattab2020colbert} & Tính MaxSim giữa các ma trận token vector & Nắm bắt tương tác token chi tiết mà vẫn tìm kiếm nhanh & Chi phí lưu trữ chỉ mục vector tăng gấp nhiều lần \\
\bottomrule
\end{tabularx}
\end{table}

\subsection{Kỹ thuật tái xếp hạng chuyên sâu bằng mô hình tương tác chéo}
Các mô hình nhúng ngữ nghĩa phổ biến như Sentence-BERT hay BGE hoạt động theo kiến trúc hai nhánh mã hóa độc lập bi-encoder: câu truy vấn và đoạn tài liệu được mã hóa thành các vector hoàn toàn tách biệt, cho phép tính toán độ tương đồng cực kỳ nhanh chóng. Tuy nhiên, do tính chất mã hóa độc lập, mô hình hai nhánh không thể mô phỏng sự tương tác tương hỗ chi tiết giữa các từ tố của câu truy vấn và các từ tố trong tài liệu.

Ngược lại, mô hình tương tác chéo cross-encoder~\cite{nogueira2019passage} tiếp nhận đồng thời cả cặp truy vấn và tài liệu cùng đi qua toàn bộ các lớp chú ý của mạng Transformer. Nhờ đó, mô hình có thể nắm bắt trọn vẹn sự tương tác ngữ nghĩa đa chiều giữa hai văn bản, mang lại độ chính xác phân loại và xếp hạng vượt trội so với kiến trúc hai nhánh, đổi lại chi phí tính toán suy luận cao hơn đáng kể. Bảng~\ref{tab:bi_vs_cross} phân tích chi tiết sự đối sánh giữa hai kiến trúc mô hình này.

\begin{table}[H]
\centering
\caption{So sánh đặc tính kỹ thuật giữa kiến trúc Bi-Encoder và Cross-Encoder (Nguồn: Tác giả tổng hợp)}
\label{tab:bi_vs_cross}
\small
\begin{tabularx}{\textwidth}{>{\raggedright\arraybackslash\bfseries}p{3.0cm} >{\raggedright\arraybackslash}X >{\raggedright\arraybackslash}X}
\toprule
\textbf{Tiêu chí so sánh} & \textbf{Kiến trúc hai nhánh Bi-Encoder} & \textbf{Kiến trúc tương tác chéo Cross-Encoder} \\
\midrule
Cơ chế mã hóa & Mã hóa độc lập truy vấn $q$ và văn bản $d$ thành hai vector tách biệt & Nối chuỗi cặp $(q, d)$ và đưa qua toàn bộ mạng Transformer đồng thời \\
Tương tác chú ý & Tự chú ý nội bộ cho từng chuỗi $q$ và $d$ riêng rẽ & Toàn bộ token của $q$ chú ý chéo đến toàn bộ token của $d$ \\
Lập chỉ mục trước & Có thể tính trước và lưu trữ toàn bộ vector tài liệu trong cơ sở dữ liệu vector & Bắt buộc tính toán trực tuyến theo thời gian thực \\
Độ phức tạp tính toán & $O(L_q^2 + L_d^2)$ cho mã hóa; $O(d)$ cho phép nhân vô hướng & $O((L_q + L_d)^2)$ cho mỗi cặp ứng viên $(q, d)$ \\
Tốc độ truy vấn & Cực nhanh (mili-giây), hỗ trợ thuật toán tìm kiếm lân cận gần đúng trên tập dữ liệu lớn & Chậm hơn đáng kể, không khả thi khi quét trực tiếp toàn bộ cơ sở dữ liệu \\
Vai trò trong hệ thống & \textbf{Pha 1: Lọc ứng viên sơ tuyển} trên diện rộng & \textbf{Pha 2: Tái xếp hạng chính xác} trên tập ứng viên chọn lọc \\
\bottomrule
\end{tabularx}
\end{table}

Chính vì vậy, trong kiến trúc truy xuất thông tin hiện đại, hai mô hình này được kết hợp theo mô hình \textbf{phễu truy xuất đa tầng}: ở tầng thứ nhất, bi-encoder và BM25 đảm nhiệm vai trò truy xuất nhanh để lọc ra một tập hợp các đoạn văn bản tương đối rộng với quy mô từ vài chục đến hàng trăm đoạn; ở tầng thứ hai, mô hình cross-encoder sẽ đóng vai trò tái xếp hạng để đánh giá lại điểm số và sắp xếp lại thứ tự ưu tiên của riêng danh sách các đoạn văn bản đó. Đây cũng là lý do vì sao khoảng cách giữa \textit{độ phủ thực tế} và \textit{trần độ phủ} của danh sách các đoạn văn bản ban đầu có ý nghĩa sống còn: nếu đoạn tài liệu chứa đáp án chính xác không lọt vào nhóm $top\text{-}k$ kết quả ở bước truy xuất sơ bộ, thì không có cơ chế tái xếp hạng nào có thể khôi phục lại được tài liệu đó ở các bước tiếp theo.


\chapter{PHƯƠNG PHÁP THỰC HIỆN VÀ KIẾN TRÚC HỆ THỐNG}
\label{chapter:phuong_phap_thuc_hien}

Chương này trình bày chi tiết về kiến trúc tổng thể, quy trình công nghệ và các giải pháp kỹ thuật được nghiên cứu, triển khai nhằm xây dựng hệ thống trợ lý ảo thông minh hỗ trợ dạy và học môn Khoa học tự nhiên cấp Trung học cơ sở. Nội dung trọng tâm bao gồm kiến trúc hệ thống Client--Server, quy trình trích xuất và tiền xử lý dữ liệu đa phương thức (ETL Pipeline), giải pháp xử lý chuyên biệt cho công thức khoa học, luồng xử lý truy vấn tăng cường truy xuất (RAG Pipeline) kết hợp truy xuất lai và tái xếp hạng, cùng thiết kế giao diện tương tác người dùng hiện đại.

\section{Kiến trúc tổng thể hệ thống}
Hệ thống trợ lý ảo được xây dựng dựa trên mô hình phân tầng Client--Server hiện đại, đảm bảo tính tách biệt giữa logic xử lý nghiệp vụ và giao diện tương tác, bao gồm hai phân hệ chính: Phân hệ máy chủ xử lý (Backend) và Phân hệ giao diện người dùng (Frontend).

\begin{figure}[htbp]
    \centering
    \includegraphics[width=0.9\textwidth]{src/images/chapter3/architecture.png}
    \caption{Sơ đồ kiến trúc tổng thể của hệ thống Trợ lý ảo RAG đa phương thức}
    \label{fig:system_architecture}
\end{figure}

Phân hệ Backend được hiện thực hoá bằng ngôn ngữ Python, sử dụng Flask làm nền tảng cung cấp dịch vụ giao diện lập trình ứng dụng (API Server), kết hợp cùng framework LangChain~\cite{langchain2022} đóng vai trò điều phối (orchestration) toàn diện các tác vụ xử lý ngữ cảnh, truy vấn cơ sở dữ liệu vector và triệu gọi mô hình ngôn ngữ lớn. Phân hệ Frontend được xây dựng trên nền tảng Next.js, mang lại giao diện ứng dụng một trang (Single Page Application --- SPA) mượt mà, phản hồi tức thời và tối ưu trải nghiệm người dùng.

\section{Quy trình thu thập và xử lý dữ liệu (ETL Pipeline)}
Quy trình Trích xuất, Biến đổi và Nạp dữ liệu (Extract, Transform, Load --- ETL) đóng vai trò nền tảng trong việc thiết lập Cơ sở tri thức (Knowledge Base) chuẩn mực cho hệ thống trợ lý ảo. Khác với các kho học liệu thuần văn bản, sách giáo khoa môn Khoa học tự nhiên sở hữu cấu trúc thông tin đa dạng và phức tạp: các đoạn văn lý thuyết đan xen với hệ thống \textbf{công thức định lượng} (Vật lí, Hoá học) và tài nguyên trực quan đa dạng (sơ đồ tế bào, sơ đồ giải phẫu, mạch điện, đồ thị động học, bảng biểu số liệu...). Nhằm bảo toàn trọn vẹn đặc trưng tri thức từ 12 cuốn sách giáo khoa thuộc ba bộ sách chuẩn mực hiện hành~\cite{sgk_kntt6, sgk_kntt7, sgk_kntt8, sgk_kntt9, sgk_ctst6, sgk_ctst7, sgk_ctst8, sgk_ctst9, sgk_cd6, sgk_cd7, sgk_cd8, sgk_cd9}, quy trình ETL được tổ chức thành hai luồng xử lý độc lập nhưng đồng bộ về mặt siêu dữ liệu.

\begin{table}[H]
\centering
\caption{Đặc trưng xuất bản và cấu trúc dữ liệu của ba bộ sách giáo khoa}
\label{tab:sgk_features}
\small
\begin{tabular}{lp{3.5cm}p{3.2cm}p{4cm}}
\toprule
\textbf{Bộ sách giáo khoa} & \textbf{Nhà xuất bản} & \textbf{Độ phân giải chuẩn} & \textbf{Đặc trưng bố cục \& Nhãn hình} \\
\midrule
\textbf{Kết nối tri thức} & NXB Giáo dục Việt Nam & $1094 \times 1536$ điểm ảnh & Nhãn hình chữ trắng trên nền khối màu bo góc; bố cục 1 cột kết hợp khối hoạt động khám phá \\
\textbf{Chân trời sáng tạo} & NXB Giáo dục Việt Nam & $2280 \times 3201$ điểm ảnh & Nhãn hình chữ đen in thường; mật độ sơ đồ thí nghiệm và chu trình hóa sinh cao \\
\textbf{Cánh Diều} & NXB Đại học Sư phạm & $2280 \times 3201$ điểm ảnh & Nhãn hình chữ đen in nghiêng; trình bày đa cột linh hoạt, mật độ đồ thị vật lý dày \\
\bottomrule
\end{tabular}
\end{table}

\subsection{Quy trình tiền xử lý dữ liệu văn bản (Text ETL)}
Luồng Text ETL chịu trách nhiệm nhận diện cấu trúc, trích xuất nội dung ký tự từ các trang sách giáo khoa và nạp vào cơ sở dữ liệu vector ChromaDB.

\textbf{Chuẩn hoá dữ liệu nguồn dạng tệp ảnh độ phân giải gốc:} Toàn bộ 12 cuốn sách giáo khoa được tổ chức và lưu trữ dưới dạng ảnh quét định dạng PNG (mỗi trang là một tệp ảnh độc lập) với độ phân giải nguyên bản từ nhà xuất bản (Bảng~\ref{tab:sgk_features}). Phương pháp này triệt tiêu hoàn toàn sự phụ thuộc vào các tham số ẩn khi kết xuất trung gian từ tệp PDF (như thư viện poppler), qua đó ngăn ngừa nguy cơ phải tái hiệu chuẩn các ngưỡng phân ngưỡng hình học của thuật toán phân tích bố cục mỗi khi thay đổi tham số DPI. Do ba bộ sách có sự khác biệt rõ rệt về độ phân giải chuẩn --- bộ Kết nối tri thức đạt $1094 \times 1536$ điểm ảnh, trong khi Chân trời sáng tạo và Cánh Diều đạt $2280 \times 3201$ điểm ảnh (riêng bộ Cánh Diều lớp 6 là $2480 \times 3480$ điểm ảnh) --- nên mọi ngưỡng tính toán kích thước hộp bao hình học đều được chuẩn hóa theo tỷ lệ tương đối so với chiều rộng trang thay vì sử dụng các giá trị điểm ảnh cố định.

\textbf{Xác thực số trang in chính xác thông qua tệp kê khai (Manifest):} Nhằm đảm bảo mọi trích dẫn đều có thể truy nguyên chính xác về trang sách vật lý mà học sinh đang sử dụng, mỗi cuốn sách được thiết lập một tệp kê khai ánh xạ tường minh giữa chỉ số tệp ảnh và \textit{số trang in thực tế}. Tệp manifest này được xây dựng dựa trên việc nhận dạng số trang in tại các góc trang kết hợp điều kiện ràng buộc tính đơn điệu tăng dần; các trang bìa được định danh riêng biệt và loại bỏ khỏi chỉ mục tri thức. Hệ thống chủ động dừng quy trình và đưa ra cảnh báo nếu phát hiện thiếu hụt tệp manifest, kiên quyết loại bỏ cơ chế gán nhãn suy đoán dựa trên chỉ số tệp vật lý nhằm tránh dẫn nguồn sai lệch cho người học.

Quy trình xử lý văn bản chi tiết gồm 4 bước chính:
\begin{enumerate}
    \item \textbf{Phân vùng bố cục và Nhận dạng ký tự quang học (OCR):} Thay vì áp dụng nhận dạng ký tự trên toàn bộ trang tài liệu một cách thô sơ (dễ gây xáo trộn thứ tự đọc giữa các cột văn bản, hộp thông tin phụ và hình ảnh), hệ thống tích hợp các kỹ thuật thị giác máy tính cổ điển để \textbf{phân vùng bố cục trang} (layout segmentation) thành các khối nội dung độc lập (hộp nội dung chính, khung hoạt động khám phá, thanh thông tin bên lề). Tiếp đó, thư viện Tesseract OCR~\cite{smith2007tesseract} với gói ngôn ngữ tiếng Việt (\texttt{tesseract-ocr-vie}) được áp dụng riêng trên từng vùng với chế độ phân đoạn trang (Page Segmentation Mode --- PSM) tương thích. Kết quả thực nghiệm kiểm chứng cho thấy phương pháp phân vùng trước khi OCR giúp giảm $3{,}8\%$ tỷ lệ mất mát từ ngữ so với phương pháp nhận dạng thô toàn trang.
    \item \textbf{Phân mảnh văn bản (Text Chunking):} Do giới hạn cửa sổ ngữ cảnh của các mô hình ngôn ngữ và nhằm tối ưu hoá độ chính xác của bước truy xuất, văn bản sau OCR được phân đoạn bằng công cụ \texttt{RecursiveCharacterTextSplitter} của LangChain với kích thước đoạn \texttt{chunk\_size = 400} ký tự và độ chồng lấn \texttt{chunk\_overlap = 120} ký tự. Kích thước 400 ký tự được xác định nhằm bao hàm trọn vẹn một luận điểm hoặc định nghĩa khoa học hoàn chỉnh, trong khi khoảng chồng lấn 120 ký tự đảm bảo tính liên kết ngữ nghĩa liền mạch giữa các đoạn tiếp giáp. Đặc biệt, bộ phân đoạn tuân thủ nguyên tắc giới hạn phân chia \textit{bên trong từng trang tài liệu}, do đó \textbf{tuyệt đối không có phân đoạn nào vắt ngang qua hai trang sách khác nhau}.
    \item \textbf{Nhúng ngữ nghĩa và Lập chỉ mục Vector:} Toàn bộ các đoạn văn bản được chuyển đổi thành các vector đặc trưng $1024$ chiều thông qua mô hình \texttt{BAAI/bge-m3}~\cite{chen2024bgem3} và lưu trữ trong collection \texttt{biology\_text} của ChromaDB. Mô hình này sở hữu khả năng nắm bắt ngữ cảnh dài và thấu hiểu đa ngữ vượt trội so với các kiến trúc nhúng nông ($384$ chiều) trên thang đo MTEB~\cite{muennighoff2022mteb}. Toàn bộ không gian vector được chuẩn hóa đồng nhất ở kích thước $1024$ chiều nhằm đảm bảo tính chuẩn xác cho các phép toán khoảng cách cosine trong cơ sở dữ liệu vector.
    \item \textbf{Khởi tạo chỉ mục thưa BM25 đồng bộ:} Song song với chỉ mục ngữ nghĩa theo vector đặc (Dense Index), hệ thống xây dựng một chỉ mục thưa Okapi BM25~\cite{robertson2009bm25} trên chính tập hợp các đoạn văn bản đó và chia sẻ chung định danh phân đoạn (chunk ID), sẵn sàng phục vụ kiến trúc truy xuất lai trình bày tại Mục~\ref{sec:relevance_gate}.
\end{enumerate}

\subsection{Giải pháp xử lý chuyên biệt cho công thức Hoá học và Vật lí}
\label{sec:formula_hybrid}
Một thách thức kỹ thuật đặc thù trong số hoá học liệu Khoa học tự nhiên là sự suy thoái nghiêm trọng của các công thức định lượng qua các công cụ OCR thông thường. Thực nghiệm khảo sát cho thấy mặc dù Tesseract nhận dạng tốt văn xuôi tiếng Việt, công cụ này \textbf{hầu như phá vỡ hoàn toàn các chỉ số dưới trong công thức hoá học và vật lí}: ví dụ phân tử \ce{O2} thường bị nhận dạng nhầm thành chuỗi \texttt{0,}, và \ce{CO2} biến dạng thành \texttt{CO,} hoặc \texttt{(0,}. Thống kê định lượng trên đường cơ sở nhận dạng ký tự truyền thống (baseline Tesseract OCR đơn thuần) ghi nhận tỷ lệ biến dạng so với nhận dạng đúng lên tới \textbf{256:3} ở bộ Cánh Diều, \textbf{377:3} ở Chân trời sáng tạo và \textbf{408:4} ở Kết nối tri thức; đồng thời không có bất kỳ ký tự chỉ số dưới dạng Unicode chuẩn nào tồn tại sau quá trình nhận dạng. Phân tích thực nghiệm cũng bác bỏ giả thuyết cho rằng sai số này do độ phân giải thấp của ảnh nguồn, bởi ngay cả các bộ sách có độ phân giải cao gấp đôi vẫn chịu sự suy giảm tương đương về tỷ lệ công thức lỗi.

\textbf{Lựa chọn công cụ thông qua thực nghiệm đối sánh độc lập (Bake-off):} Nhằm tìm kiếm giải pháp thay thế hiệu quả, đề tài thiết lập một tập kiểm thử chuẩn (gold set) gồm $97$ ô văn bản phân bố trên $15$ trang sách (cân đối $5$ trang cho mỗi bộ sách) được các chuyên gia đối chiếu và nhập liệu thủ công. Kết quả thử nghiệm mô hình thị giác--ngôn ngữ tiên tiến MinerU2.5 cho thấy mô hình này cải thiện vượt bậc khả năng đọc chỉ số dưới với tỷ lệ chính xác đạt $\mathbf{0{,}441}$ (cao gấp $9{,}2$ lần so với mức $0{,}048$ của Tesseract). Tuy nhiên, MinerU2.5 lại bộc lộ nhược điểm lớn ở văn xuôi tiếng Việt khi làm tăng tỷ lệ lỗi dấu thanh lên mức $\mathbf{0{,}037}$, cao hơn đáng kể so với mức $0{,}016$ của Tesseract. Do nội dung văn xuôi chiếm hơn $93\%$ dung lượng sách giáo khoa, việc thay thế hoàn toàn Tesseract bằng MinerU2.5 sẽ dẫn tới sự suy giảm nghiêm trọng về chất lượng ngữ nghĩa tiếng Việt của toàn bộ tài liệu.

\textbf{Kiến trúc xử lý lai phân cấp theo dòng (Line-level Hybrid Architecture):} Từ phát hiện trên, đề tài đề xuất giải pháp kiến trúc lai: duy trì Tesseract làm công cụ chủ đạo cho toàn bộ văn bản thông thường, đồng thời phát triển một cổng phân loại nhanh (\texttt{is\_formula\_suspect}) để quét từng dòng văn bản và phát hiện các mẫu nghi ngờ chứa công thức bị phân rã. Bộ lọc này đạt độ chính xác (Precision) $\mathbf{0{,}8654}$ và độ phủ (Recall) $\mathbf{1{,}0000}$ trên tập kiểm định gồm $89$ ô văn bản. Chỉ những dòng bị nghi ngờ lỗi mới được trích xuất ảnh tương ứng để gửi tới mô hình MinerU2.5 xử lý lại. 

Nhằm đảm bảo tính nguyên vẹn và tránh việc hệ thống tự ý chắp vá thông tin sai lệch, cơ chế ghép ngược (merging) tuân thủ nguyên tắc ràng buộc chặt chẽ: phép thay thế chỉ được kích hoạt khi số lượng thực thể khoa học do MinerU2.5 phát hiện trùng khớp hoàn toàn với số lượng vị trí suy biến do Tesseract để lại. Trường hợp không thỏa mãn điều kiện khớp số lượng, dòng văn bản sẽ được bảo lưu nguyên trạng kết quả của Tesseract. Quy trình này đã được triển khai xử lý hoàn tất trên toàn bộ $16.515$ phân đoạn của kho tri thức thực tế (ghi nhận $4\,293$ lượt ghép dòng thành công, $0$ lỗi ngoại lệ).

\subsection{Quy trình tiền xử lý dữ liệu hình ảnh (Image ETL)}
Nhằm hỗ trợ học sinh tiếp cận kiến thức thông qua các tài nguyên trực quan, luồng Image ETL được tổ chức tuần tự theo 4 giai đoạn:
\begin{enumerate}
    \item \textbf{Phát hiện khung ảnh dựa trên neo chú thích (Anchor-first Detection):} Thay vì phụ thuộc hoàn toàn vào các mô hình học sâu phát hiện vật thể (vốn dễ chịu ảnh hưởng bởi nhiễu nền của trang sách phức tạp), hệ thống ưu tiên định vị các chuỗi ký tự định danh chính thức có cấu trúc \texttt{Hình A.B} xuất hiện trên trang. Dựa trên quy ước bố cục in ấn của từng bộ sách, toạ độ của hình ảnh minh hoạ phía trên hoặc phía dưới dòng chú thích sẽ được trích xuất chính xác. Mô hình thị giác \texttt{google/owlvit-base-patch32}~\cite{minderer2022owlvit} chỉ đóng vai trò hỗ trợ phụ trợ đối với các khung ảnh đặc thù không đi kèm nhãn chú thích tiêu chuẩn. Ưu thế then chốt của giải pháp này là tính minh bạch và khả năng giải trình: mỗi khung ảnh được cắt ra đều gắn liền với một nhãn định danh chính thống trên trang sách.
    \item \textbf{Cấu hình trích xuất thích ứng theo định dạng xuất bản:} Thống kê bố cục cho thấy các bộ sách áp dụng phong cách thiết kế nhãn hình khác nhau: bộ Kết nối tri thức sử dụng nhãn \textbf{chữ trắng in trên nền khối màu bo góc} (phát hiện trung bình $23$--$28$ nhãn/30 trang qua kênh xử lý ngưỡng màu chuyên biệt); trong khi Chân trời sáng tạo và Cánh Diều sử dụng \textbf{chữ đen thông thường} ($30$--$65$ nhãn/30 trang qua kênh OCR chuẩn). Hệ thống tự động nhận diện hồ sơ bố cục (layout profile) tương ứng với từng bộ sách để áp dụng giải thuật bóc tách tối ưu, ngăn ngừa hoàn toàn hiện tượng nhầm lẫn cấu trúc.
    \item \textbf{Xây dựng siêu dữ liệu xác thực từ điểm ảnh:} Mỗi khung hình sau khi trích xuất được gắn kèm các trường siêu dữ liệu trích xuất trực tiếp từ trang sách gốc: nhãn hình (\texttt{Hình A.B}), chuỗi chú thích nội dung và các từ ngữ nhận dạng được bên trong khung ảnh. Hệ thống chủ động loại bỏ việc sử dụng các mô hình tự động sinh chú thích ảnh (Image Captioning) như Vintern-1B. Khảo nghiệm trên dữ liệu thực tế cho thấy mô hình sinh chú thích tiêu tốn tới $17{,}6$ giây/ảnh trên CPU, đồng thời bộc lộ tỷ lệ sinh ảo giác cao ($4/12$ trường hợp tự suy diễn chi tiết không có thật, và gán sai $4/4$ số hiệu hình). Việc sử dụng siêu dữ liệu xác thực từ văn bản gốc đảm bảo độ chính xác tuyệt đối và loại bỏ hoàn toàn các rủi ro sai lệch tri thức trực quan.
    \item \textbf{Lập chỉ mục đa phương thức (Multimodal Indexing):} Mỗi khung hình được mã hoá thành vector nhúng $512$ chiều thông qua mô hình liên kết ngôn ngữ -- thị giác \texttt{clip-vit-base-patch16}~\cite{radford2021clip} và lưu trữ trong collection \texttt{biology\_images} cùng siêu dữ liệu định danh, phục vụ việc tìm kiếm tương đồng đa phương thức.
\end{enumerate}

Bên cạnh quy trình tự động, hệ thống cung cấp sẵn các công cụ kết xuất dữ liệu phục vụ cơ chế chuyên gia rà soát (Human-in-the-loop), tạo điều kiện cho các chuyên gia chuyên môn hoặc người quản trị hệ thống có thể kiểm tra và hiệu chỉnh siêu dữ liệu hình ảnh khi có nhu cầu chuẩn hóa chuyên sâu.

\section{Luồng truy xuất và sinh văn bản (RAG Pipeline)}
Khi người dùng (học sinh hoặc giáo viên) nhập câu hỏi trên giao diện, hệ thống Backend tiếp nhận và điều phối luồng xử lý RAG đa phương thức theo các giai đoạn được mô tả trong Hình~\ref{fig:rag_pipeline_detail}:

\begin{figure}[htbp]
    \centering
    \includegraphics[width=0.65\textwidth]{src/images/chapter3/rag_pipeline.png}
    \caption{Luồng xử lý truy vấn RAG đa phương thức tại Backend}
    \label{fig:rag_pipeline_detail}
\end{figure}

\begin{enumerate}
    \item \textbf{Cơ chế truy xuất kép (Dual Retrieval):} Phía dữ liệu văn bản, câu hỏi được chuyển đồng thời qua kênh tìm kiếm ngữ nghĩa vector và kênh tìm kiếm từ khoá BM25 (mỗi kênh trích chọn $20$ ứng viên ban đầu). Tập hợp ứng viên sau đó được hợp nhất và tái xếp hạng chuyên sâu bằng mô hình cross-encoder để chọn ra $\mathtt{RETRIEVER\_K = 3}$ đoạn văn bản có độ liên quan cao nhất. Phía dữ liệu hình ảnh, hệ thống trích chọn $\mathtt{IMAGE\_RETRIEVER\_K = 5}$ hình ảnh tiêu biểu dựa trên sự kết hợp giữa độ tương đồng vector CLIP và phép so khớp chuỗi siêu dữ liệu có dấu.
    \item \textbf{Kỹ nghệ thiết kế lời nhắc (Prompt Engineering):} Ngữ cảnh văn bản được chọn lọc sẽ được liên kết đồng bộ với thông tin mô tả của các hình ảnh tìm được. Lời nhắc (Prompt) được thiết kế theo cấu trúc sư phạm chuyên biệt, định hướng mô hình ngôn ngữ lớn đóng vai trò người trợ giảng Khoa học tự nhiên chuẩn mực, trả lời gãy gọn, mạch lạc và tuyệt đối trung thành với ngữ cảnh tài liệu được cung cấp.
    \item \textbf{Sinh câu trả lời (Text Generation):} Mô hình Qwen2.5-3B-Instruct~\cite{qwen25_2024} tiếp nhận lời nhắc hoàn chỉnh để sinh ra câu trả lời cuối cùng. Siêu tham số nhiệt độ được thiết lập ở mức thấp (\texttt{temperature = 0.1}) nhằm tối đa hóa tính xác thực của câu trả lời, hạn chế tối đa tính ngẫu nhiên và ngăn ngừa hiện tượng suy diễn sai lệch kiến thức khoa học.
\end{enumerate}

\subsection{Cơ chế định tuyến theo ý định truy vấn (Query Intent Routing)}
Không phải mọi tương tác của học sinh đều đòi hỏi đồng thời cả văn bản lý thuyết lẫn tài liệu hình ảnh. Module định tuyến ý định (\texttt{query\_intent}) thực hiện phân tích cấu trúc ngữ nghĩa của câu hỏi trước khi kích hoạt quy trình truy xuất. Khi phát hiện các truy vấn thể hiện rõ ý định quan sát trực quan (chứa các cụm từ như "cho xem hình", "sơ đồ cấu tạo", "biểu đồ minh hoạ"), hệ thống sẽ ưu tiên tối đa luồng truy xuất hình ảnh. Ngược lại, đối với các câu hỏi bản chất lý thuyết trừu tượng hoặc bài tập tính toán, luồng văn bản sẽ đóng vai trò chủ đạo nhằm loại bỏ các hình ảnh gây nhiễu ngữ cảnh.

\subsection{Truy xuất lai thưa -- dày và cơ chế tái xếp hạng}
\label{sec:relevance_gate}
Truy xuất ngữ nghĩa vector (Dense Retrieval) sở hữu khả năng bao quát ý nghĩa tương đồng xuất sắc, song lại gặp bất lợi lớn trước các yêu cầu so khớp định danh chính xác vốn rất phổ biến trong môn Khoa học tự nhiên: ví dụ truy vấn về phân tử \ce{CO2} hoặc định luật Boyle--Mariotte đòi hỏi sự xuất hiện chuẩn xác của chuỗi ký hiệu thay vì xấp xỉ ngữ nghĩa. Ngược lại, truy xuất từ khóa thưa (Sparse Retrieval) lại mất khả năng liên kết khi học sinh sử dụng ngôn ngữ diễn đạt khác biệt với nguyên văn sách giáo khoa. Hệ thống giải quyết sự đánh đổi này thông qua kiến trúc \textbf{Truy xuất lai (Hybrid Retrieval)}: chạy song song kênh vector \texttt{bge-m3} và kênh Okapi BM25~\cite{robertson2009bm25} trên cùng cơ sở dữ liệu phân đoạn, sau đó tiến hành hợp nhất kết quả.

\subsubsection{Tối ưu hóa chỉ mục BM25 và chuẩn hoá công thức}
Thuật toán Okapi BM25 được tự hiện thực hoá trong hệ thống nhằm hỗ trợ việc tinh chỉnh động hai siêu tham số $k_1$ và $b$ ngay tại thời điểm xử lý truy vấn. Quá trình quét lưới thực nghiệm $5 \times 5$ đã xác định bộ tham số tối ưu đạt $k_1 = 0{,}7$ và $b = 0{,}75$.

Một phát hiện thực nghiệm có ý nghĩa khoa học quan trọng là: thuật toán tách từ \textbf{bảo lưu dấu thanh tiếng Việt} đạt hiệu năng vượt trội so với giải pháp loại bỏ dấu thanh (chỉ số MRR đạt $0{,}820$ so với $0{,}755$). Việc loại bỏ dấu thanh dẫn tới hiện tượng đồng âm sai lệch nghiêm trọng giữa các hư từ và các khái niệm khoa học then chốt (ví dụ: \textit{khí} biến thành \textit{khi}, \textit{độ} thành \textit{do}, \textit{lá} thành \textit{la}, \textit{đá} thành \textit{da}).

Đối với các công thức khoa học, hệ thống áp dụng cơ chế \textbf{chuẩn hoá ký hiệu ở phía truy vấn và phía chỉ mục thưa, tuyệt đối không chỉnh sửa nội dung văn bản lưu trữ gốc}. Thử nghiệm trên $12$ công thức hóa học tiêu biểu cho thấy số phân đoạn chính xác lọt vào top-10 tăng từ $6$ lên $97$ đoạn, và số câu truy vấn tìm được tài liệu đích tăng ngoạn mục từ $1/12$ lên $11/12$.

\subsubsection{Quy trình hợp nhất và Tái xếp hạng bằng Cross-Encoder}
Quy trình truy xuất được tổ chức tuần tự: \textbf{Truy xuất song song $\rightarrow$ Hợp nhất danh sách $\rightarrow$ Tái xếp hạng}. Bước tái xếp hạng sử dụng mô hình cross-encoder \texttt{BAAI/bge-reranker-v2-m3} cùng họ BGE, thực hiện tính toán lại điểm số liên quan trên toàn bộ tập ứng viên hợp nhất từ hai kênh.

Các phân tích thực nghiệm tại Chương~\ref{chapter:thuc_nghiem_danh_gia} đã chứng minh rằng việc áp dụng thêm cổng lọc dựa trên khoảng cách tương đối trước bước tái xếp hạng không mang lại cải thiện, thậm chí làm suy giảm độ phủ (Recall@10 giảm từ $1{,}000$ xuống $0{,}890$) do loại bỏ nhầm các đoạn tài liệu đúng ở các câu hỏi phức hợp nhiều ý. Thay vào đó, việc thiết lập \textbf{ngưỡng điểm số tuyệt đối} tại bộ tái xếp hạng mang lại hiệu quả phân loại tin cậy hơn hẳn. Do đó, cấu hình tối ưu mặc định của hệ thống được xác lập là: \textbf{Truy xuất lai, kích hoạt tái xếp hạng bằng cross-encoder và vô hiệu hoá cổng lọc khoảng cách tương đối}.

\begin{table}[H]
\centering
\caption{Tổng hợp siêu tham số vận hành của hệ thống trợ lý ảo}
\label{tab:hyperparams}
\small
\begin{tabular}{lp{6.2cm}c}
\toprule
\textbf{Thành phần} & \textbf{Ý nghĩa kỹ thuật} & \textbf{Giá trị cấu hình} \\
\midrule
\texttt{EMBEDDING\_DIM} & Số chiều không gian vector BGE-M3 & $1024$ \\
\texttt{CLIP\_DIM} & Số chiều không gian vector hình ảnh CLIP & $512$ \\
\texttt{CHUNK\_SIZE} & Kích thước phân đoạn văn bản & $400$ ký tự \\
\texttt{CHUNK\_OVERLAP} & Khoảng trượt liên kết giữa các đoạn liền kề & $120$ ký tự \\
\texttt{BM25\_K1} & Tham số bão hòa tần suất từ khóa BM25 & $0{,}7$ \\
\texttt{BM25\_B} & Tham số chuẩn hóa độ dài văn bản BM25 & $0{,}75$ \\
\texttt{TOP\_K\_CANDIDATES} & Số ứng viên sơ tuyển mỗi kênh (BM25, Dense) & $20$ \\
\texttt{RETRIEVER\_K} & Số đoạn văn bản chọn lọc sau tái xếp hạng & $3$ \\
\texttt{IMAGE\_RETRIEVER\_K} & Số khung hình trích chọn đa phương thức & $5$ \\
\texttt{RERANK\_SCORE\_MIN} & Ngưỡng điểm tối thiểu của Cross-Encoder & $0{,}59$ \\
\texttt{LLM\_TEMPERATURE} & Nhiệt độ sinh văn bản của Qwen2.5 & $0{,}1$ \\
\texttt{RANDOM\_SEED} & Hạt giống ngẫu nhiên tái lập thực nghiệm & $42$ \\
\bottomrule
\end{tabular}
\end{table}

\section{Thiết kế giao diện người dùng tương tác (Frontend)}
Giao diện ứng dụng được thiết kế tối ưu hóa cho trải nghiệm học tập của học sinh Trung học cơ sở, phát triển trên nền tảng \textbf{Next.js 16.2.6} và thư viện \textbf{React 19.2.4}.

Các công nghệ giao diện chính bao gồm:
\begin{itemize}
    \item \textbf{Tailwind CSS:} Đảm bảo thiết kế giao diện hiện đại, trực quan, thân thiện với lứa tuổi học đường và tương thích hoàn hảo trên đa dạng thiết bị (Responsive Design).
    \item \textbf{Framer Motion:} Hiện thực hoá các chuyển động vi mô (micro-animations) mượt mà khi người dùng tương tác gửi nhận câu hỏi, tạo cảm giác thân thiện như đang trao đổi trực tiếp cùng một gia sư chuyên môn.
    \item \textbf{Hiển thị công thức khoa học bằng KaTeX:} Tích hợp thư viện KaTeX kết hợp cơ chế tự động phân tích và chuyển đổi các ký tự hóa học phẳng (như \texttt{H2SO4}) sang định dạng ký hiệu chuẩn khoa học (\ce{H2SO4}). Quá trình này tra cứu bảng tuần hoàn $118$ nguyên tố hoá học để đảm bảo chỉ chuyển đổi khi chuỗi ký hiệu hợp lệ, tránh biến đổi nhầm các từ ngữ văn bản thông thường.
    \item \textbf{Minh bạch hóa học liệu qua khối trích dẫn chuyên biệt:} Hệ thống máy chủ Backend phân tách tường minh giữa nội dung văn bản câu trả lời (\texttt{answer\_text}) và trường siêu dữ liệu trích dẫn nguồn có cấu trúc (\texttt{citations} gồm tên bộ sách, số trang và mục bài học). Giao diện người dùng hiển thị riêng biệt khối trích dẫn này ngay bên dưới câu trả lời, giúp học sinh và giáo viên dễ dàng đối chiếu trực tiếp với trang sách giáo khoa tương ứng.
    \item \textbf{Cơ chế phục vụ tài nguyên hình ảnh trực tiếp từ máy chủ:} Thay vì sao chép thủ công kho dữ liệu hình ảnh sang phía máy khách, hệ thống xây dựng tuyến đường dẫn (route) phục vụ hình ảnh động từ máy chủ Backend. Giải pháp này đảm bảo tính đồng bộ dữ liệu hình ảnh và tối ưu hóa hiệu năng truyền tải trên mạng Internet.
\end{itemize}

\subsection{Một số giao diện tiêu biểu của ứng dụng}
Hình~\ref{fig:fe_overview} minh họa các màn hình chức năng chính trong luồng trải nghiệm của người dùng: từ quá trình đăng nhập hệ thống, khung đối thoại đặt câu hỏi cho tới màn hình hiển thị câu trả lời tích hợp công thức KaTeX và hình ảnh trích xuất. Hình~\ref{fig:fe_image_detail} thể hiện tính năng chi tiết của hệ thống đa phương thức: cho phép người dùng phóng to hình ảnh minh họa và kiểm tra trực tiếp trang sách giáo khoa gốc do hệ thống trích xuất.

\begin{figure}[H]
    \centering
    \begin{subfigure}[t]{0.48\textwidth}
        \centering
        \includegraphics[width=\linewidth]{src/images/chapter3/fe-login.png}
        \caption{Màn hình đăng nhập}
        \label{fig:fe_login}
    \end{subfigure}\hfill
    \begin{subfigure}[t]{0.48\textwidth}
        \centering
        \includegraphics[width=\linewidth]{src/images/chapter3/fe-ask-question.png}
        \caption{Màn hình đặt câu hỏi}
        \label{fig:fe_ask}
    \end{subfigure}

    \vspace{1em}
    \begin{subfigure}[t]{0.48\textwidth}
        \centering
        \includegraphics[width=\linewidth]{src/images/chapter3/fe-chatbox.png}
        \caption{Khung tương tác trò chuyện}
        \label{fig:fe_chatbox}
    \end{subfigure}\hfill
    \begin{subfigure}[t]{0.48\textwidth}
        \centering
        \includegraphics[width=\linewidth]{src/images/chapter3/fe-response.png}
        \caption{Hiển thị câu trả lời và trích dẫn}
        \label{fig:fe_response}
    \end{subfigure}
    \caption{Các màn hình chính của giao diện trợ lý ảo Khoa học tự nhiên.}
    \label{fig:fe_overview}
\end{figure}

\begin{figure}[H]
    \centering
    \includegraphics[width=0.75\textwidth]{src/images/chapter3/fe-image-detail.png}
    \caption{Chi tiết hình ảnh minh hoạ kèm trang sách gốc do hệ thống truy xuất.}
    \label{fig:fe_image_detail}
\end{figure}

\chapter{THỰC NGHIỆM, ĐÁNH GIÁ VÀ THẢO LUẬN}
\label{chapter:thuc_nghiem_danh_gia}

Chương này trình bày toàn diện quá trình thực nghiệm, kiểm thử định lượng và thảo luận khoa học về hiệu năng của hệ thống trợ lý ảo. Nội dung bao gồm: thiết lập môi trường và công cụ thực nghiệm; quy mô và cấu trúc kho tri thức số hóa; phương pháp xây dựng bộ dữ liệu kiểm thử có gắn nhãn nguồn xác thực; hệ thống các độ đo đánh giá độc lập cho hai pha truy xuất tri thức và sinh ngôn ngữ; kết quả thực nghiệm chi tiết trên \textbf{240 câu hỏi kiểm thử} phủ rộng 12 cuốn sách giáo khoa; và phần thảo luận chuyên sâu về các phát hiện khoa học, ưu điểm cùng các giới hạn kỹ thuật được nhận diện qua thực nghiệm.

\section{Môi trường và công cụ thực nghiệm}
Quá trình tiền xử lý học liệu bao gồm các công đoạn nhận dạng ký tự quang học, phân vùng bố cục, phát hiện và bóc tách hình ảnh, kết hợp với quá trình suy luận của mô hình ngôn ngữ lớn đòi hỏi cấu hình tài nguyên tính toán phù hợp. Bảng~\ref{tab:env} tổng hợp chi tiết cấu hình phần cứng và hệ thống nền tảng công nghệ thực nghiệm được sử dụng trong toàn bộ các đợt thực nghiệm của đề tài.

\begin{table}[H]
\centering
\caption{Cấu hình môi trường và hệ thống nền tảng công nghệ thực nghiệm (Nguồn: Thực nghiệm của đồ án)}
\label{tab:env}
\small
\begin{adjustbox}{max width=\textwidth}
\begin{tabular}{p{4.8cm}p{10.2cm}}
\toprule
\textbf{Thành phần} & \textbf{Chi tiết cấu hình} \\
\midrule
Nền tảng tính toán & Máy trạm cục bộ, hệ điều hành Windows 11 \\
Bộ vi xử lý và bộ nhớ chính & 16 lõi xử lý, 63,9 GB RAM \\
Bộ xử lý đồ họa GPU & \textbf{Không trang bị} (\texttt{torch 2.11.0+cpu}, CUDA không khả dụng) \\
Ngôn ngữ lập trình & Python 3.12.8 \\
Nền tảng phần mềm cốt lõi & LangChain 1.2, Transformers 4.46.3, PyTorch 2.11 (CPU) \\
Cơ sở dữ liệu véc-tơ & ChromaDB 1.5.9 \\
Mô hình sinh ngôn ngữ & Qwen2.5-3B-Instruct \\
Mô hình nhúng văn bản & \texttt{BAAI/bge-m3} (không gian nhúng 1024 chiều) \\
Chỉ mục thưa & Okapi BM25 tối ưu hóa ($k_1 = 0{,}7$, $b = 0{,}75$) \\
Mô hình tái xếp hạng & \texttt{BAAI/bge-reranker-v2-m3} (mô hình tương tác chéo) \\
Mô hình thị giác máy tính & CLIP-ViT-B/16, OwlViT-B/32 \\
Công cụ nhận dạng ký tự & Tesseract OCR tích hợp gói dữ liệu tiếng Việt \\
Mô hình giám khảo độc lập & Tập hợp 4 mô hình ngôn ngữ lớn luân phiên qua Groq API (Mục~\ref{sec:judge_pha2}) \\
\bottomrule
\end{tabular}
\end{adjustbox}
\end{table}

Một ưu thế kỹ thuật nổi bật của cấu hình thực nghiệm là \textbf{toàn bộ quy trình tiền xử lý học liệu được thiết kế vận hành độc lập trên CPU mà không đòi hỏi GPU}. Đây là một kết quả có ý nghĩa ứng dụng thực tiễn cao: hệ thống được tối ưu hóa dựa trên các kỹ thuật thị giác máy tính phân tích bố cục hình học và nhận dạng ký tự xác định, giúp giai đoạn xây dựng kho tri thức quy mô lớn hoàn toàn khả thi trên các máy trạm thông thường. Về mặt chi phí thời gian thực tế: luồng xử lý văn bản hoàn tất 12 cuốn sách trong \textbf{3 giờ 20 phút} (tốc độ trung bình $5{,}0$ giây/trang) cùng $49$ phút cho bước khởi tạo tệp manifest; luồng hình ảnh hoàn thành trong \textbf{5 giờ 01 phút} cho toàn bộ kho sách. Trong hệ thống đề xuất, tài nguyên GPU chỉ thực sự cần thiết nếu muốn gia tốc tốc độ sinh câu trả lời của mô hình ngôn ngữ hoặc khi xử lý các mô hình thị giác - ngôn ngữ chuyên sâu (như bước OCR công thức qua MinerU đã được thực hiện hoàn tất trên môi trường GPU điện toán đám mây trước khi đưa vào chỉ mục triển khai, trình bày tại Mục~\ref{sec:formula_hybrid}).

Nhằm đảm bảo tính khách quan khoa học, mô hình giám khảo chấm điểm được phân tách độc lập hoàn toàn với mô hình sinh ngôn ngữ Qwen2.5, tuân thủ nghiêm ngặt phương pháp đánh giá bằng mô hình ngôn ngữ lớn làm giám khảo độc lập~\cite{zheng2023llmjudge} nhằm loại bỏ thiên kiến tự chấm điểm thiên vị của mô hình.

\section{Dữ liệu tri thức đã số hóa và lập chỉ mục}
\label{sec:corpus}
Cơ sở tri thức của hệ thống được xây dựng từ \textbf{12 cuốn sách giáo khoa Khoa học tự nhiên} các lớp 6, 7, 8 và 9 thuộc cả ba bộ sách chuẩn mực (Cánh Diều, Chân trời sáng tạo, Kết nối tri thức với cuộc sống), được lưu trữ dưới dạng tệp ảnh PNG đơn trang ở độ phân giải gốc của nhà xuất bản, với tổng quy mô \textbf{2.399 trang tài liệu} vật lý trên đĩa. Trong đó, \textbf{2.387 trang nội dung bài học} được đưa vào lập chỉ mục chính thức; phần chênh lệch $12$ trang là các trang bìa và trang thông tin xuất bản của 12 cuốn sách, được tệp manifest định danh riêng biệt và chủ động loại bỏ khỏi cơ sở tri thức nhằm tránh gây nhiễu cho bước truy xuất ngữ cảnh. Quy trình tiền xử lý hai luồng đồng bộ xử lý văn bản và hình ảnh; toàn bộ dữ liệu sau đó được lưu trữ trong ba tập hợp dữ liệu của ChromaDB, với số liệu thống kê tổng hợp tại Bảng~\ref{tab:corpus}.

\textbf{Quy mô dữ liệu văn bản:} Sau quy trình nhận dạng ký tự theo vùng bố cục và phân mảnh (đoạn $400$ ký tự, độ chồng lấn $120$ ký tự), toàn bộ tập ngữ liệu tạo ra \textbf{16.515 đoạn văn bản}. Mỗi đoạn được chuyển đổi bằng mô hình \texttt{BAAI/bge-m3}~\cite{chen2024bgem3} thành một \textbf{vector đặc trưng 1024 chiều} trong tập hợp dữ liệu \texttt{biology\_text}. Đồng thời, chỉ mục thưa Okapi BM25 được thiết lập đồng bộ trên $16.515$ phân đoạn này với không gian từ vựng đạt \textbf{20.914 từ vựng}.

\textbf{Quy mô dữ liệu hình ảnh:} Thông qua phương pháp phát hiện dựa trên neo chú thích kết hợp mô hình OwlViT phụ trợ, hệ thống trích xuất được \textbf{3.881 khung hình minh hoạ} trên cả 12 cuốn sách (bao gồm $2.172$ hình đơn, $824$ hình con phân rã, $442$ khung hoạt động khám phá, $324$ hình ghép phức hợp, $87$ khung thông tin và $32$ nhóm thiết bị thí nghiệm). Mỗi khung hình được nhúng thành một \textbf{vector 512 chiều} thông qua mô hình \texttt{clip-vit-base-patch16} lưu trữ trong tập hợp dữ liệu \texttt{biology\_images}, đi kèm bản ghi siêu dữ liệu trong \texttt{biology\_image\_metadata}.

Toàn bộ siêu dữ liệu của mỗi khung hình được xây dựng dựa trên nguyên tắc trích xuất xác thực từ trang sách gốc (bao gồm nhãn hình, dòng chú thích và chữ nhận dạng bên trong khung), loại bỏ hoàn toàn các mô tả tự sinh từ mô hình nhằm triệt tiêu nguy cơ ảo giác tri thức trực quan.

\textbf{Độ phủ nhãn hình ảnh:} Để đánh giá khách quan độ phủ của luồng trích xuất hình ảnh trên toàn bộ ngữ liệu, hệ thống đối chiếu tập hợp nhãn \texttt{Hình A.B} xuất hiện trong \textit{văn bản đã lập chỉ mục} với tập hợp nhãn của các khung hình trích xuất được. Kết quả đạt được: bộ Cánh Diều đạt từ $92\%$ đến $97\%$, Kết nối tri thức đạt từ $95\%$ đến $97\%$, và Chân trời sáng tạo đạt từ $72\%$ đến $89\%$. Cần lưu ý rằng chỉ số này đại diện cho một cận dưới ước lượng, do một số nhãn xuất hiện trong văn bản là các trích dẫn tham chiếu trong bài học thay vì chú thích ảnh; do đó, độ đo này mang giá trị so sánh tương đối giữa các bộ sách hơn là một tỷ lệ chính xác tuyệt đối.

\begin{table}[H]
\centering
\caption{Thống kê cơ sở tri thức đã số hóa và lập chỉ mục theo từng bộ sách giáo khoa (Nguồn: Đo lường thực tế từ cơ sở dữ liệu)}
\label{tab:corpus}
\small
\begin{adjustbox}{max width=\textwidth}
\begin{tabular}{llccc}
\toprule
\textbf{Khối lớp} & \textbf{Bộ sách giáo khoa} & \textbf{Trang lập chỉ mục} & \textbf{Phân đoạn văn bản} & \textbf{Vector hình ảnh} \\
\midrule
\multirow{3}{*}{Lớp 6}
 & KHTN 6 Cánh Diều            & 178 & 1.124 & 418 \\
 & KHTN 6 Chân trời sáng tạo   & 203 & 1.280 & 470 \\
 & KHTN 6 Kết nối tri thức     & 194 & 1.125 & 286 \\
\midrule
\multirow{3}{*}{Lớp 7}
 & KHTN 7 Cánh Diều            & 170 & 1.156 & 269 \\
 & KHTN 7 Chân trời sáng tạo   & 187 & 1.316 & 335 \\
 & KHTN 7 Kết nối tri thức     & 178 & 1.073 & 203 \\
\midrule
\multirow{3}{*}{Lớp 8}
 & KHTN 8 Cánh Diều            & 206 & 1.451 & 324 \\
 & KHTN 8 Chân trời sáng tạo   & 222 & 1.764 & 423 \\
 & KHTN 8 Kết nối tri thức     & 195 & 1.253 & 216 \\
\midrule
\multirow{3}{*}{Lớp 9}
 & KHTN 9 Cánh Diều            & 214 & 1.613 & 379 \\
 & KHTN 9 Chân trời sáng tạo   & 214 & 1.814 & 324 \\
 & KHTN 9 Kết nối tri thức     & 226 & 1.546 & 234 \\
\midrule
\multicolumn{2}{r}{\textbf{Tổng cộng toàn hệ thống}} & \textbf{2.387} & \textbf{16.515} & \textbf{3.881} \\
\bottomrule
\end{tabular}
\end{adjustbox}
\end{table}

Tổng thể kho tri thức bao gồm \textbf{20.396 vector} ($16.515$ vector văn bản và $3.881$ vector hình ảnh) được lập chỉ mục và quản trị trong ChromaDB. Mật độ hình ảnh trích xuất trung bình đạt từ $1{,}1$ đến $2{,}3$ khung hình/trang, phản ánh sự cải thiện rõ rệt về tính đồng nhất sau khi chuẩn hóa thuật toán xử lý bố cục.

Về mặt cấu trúc siêu dữ liệu, đối với trường số hiệu bài học (\texttt{lesson\_id}), hệ thống tuân thủ nguyên tắc bảo toàn tính chân thực: chỉ gán giá trị khi chuỗi mục lục được trích xuất liền mạch và tường minh (hiện đạt $4.919/16.515$ đoạn văn bản, tương ứng $29{,}8\%$), chủ động để trống đối với các trường hợp chưa đọc được trục mục lục thay vì áp dụng phỏng đoán ngẫu nhiên nhằm loại trừ nguy cơ lan truyền sai lệch thông tin bài học.

\section{Xây dựng bộ dữ liệu kiểm thử chuẩn mực}
\label{sec:testset}
Nhằm đảm bảo quá trình đánh giá hiệu năng diễn ra hoàn toàn khách quan, minh bạch và có khả năng tái lập, đề tài xây dựng quy trình khởi tạo dữ liệu kiểm thử gắn nhãn nguồn xác định. Đây là tiền đề bắt buộc để lượng hoá chính xác các chỉ số truy xuất thông tin chuẩn mực gồm Precision, Recall và MRR dựa trên sự thật thay vì phụ thuộc thuần túy vào cảm tính hoặc phán đoán định tính.

Phương pháp kiểm thử được thiết lập dựa trên nguyên tắc lấy mẫu \textbf{ngẫu nhiên phân tầng trên toàn bộ tập ngữ liệu}, loại bỏ các hạn chế của việc chọn mẫu thủ công cục bộ.

\subsection{Nguyên lý thiết lập nhãn chuẩn đối sánh}
Bộ dữ liệu kiểm thử bao gồm \textbf{ba nhóm câu hỏi} đặc thù:
\begin{itemize}
    \item \textbf{Nhóm câu hỏi văn bản} và \textbf{Nhóm câu hỏi hình ảnh:} Mỗi câu hỏi được khởi tạo tương ứng từ một phân đoạn văn bản hoặc một khung hình cụ thể trong chỉ mục. Do đó, tài liệu chuẩn đối sánh được xác lập chính xác thông qua cặp định danh \texttt{(tên\_sách, số\_trang)} trích xuất trực tiếp từ siêu dữ liệu gốc.
    \item \textbf{Nhóm câu hỏi ngoài phạm vi:} Nhóm này tập hợp các câu hỏi thuộc các môn học khác ngoài Khoa học tự nhiên (Lịch sử, Địa lí, Ngữ văn, Tin học, Toán học). Về mặt thiết kế, nhóm này \textbf{không có tài liệu nguồn sách giáo khoa KHTN tương ứng}, và phản hồi chuẩn của hệ thống là \textbf{phải nhận biết câu hỏi ngoài phạm vi và từ chối trả lời}, đóng vai trò kiểm chứng định lượng cho khả năng phòng chống hiện tượng ảo giác tri thức (Mục~\ref{sec:ablation}).
\end{itemize}

\subsection{Quy trình khởi tạo và kiểm định dữ liệu kiểm thử}
Quy trình khởi tạo bộ kiểm thử được tự động hóa kết hợp kiểm định chuyên gia (Thuật toán~\ref{alg:gen}), sử dụng hạt giống ngẫu nhiên cố định \texttt{seed=42} nhằm đảm bảo tính tái lập. Tỷ lệ câu hỏi hình ảnh $p_{hinh}$ được xác định động dựa trên tương quan kích thước thực tế của tập chỉ mục:
\begin{equation}
p_{hinh} = \frac{n_{anh}}{n_{chunk} + n_{anh}} \approx \frac{3\,780}{11\,444 + 3\,780} \approx 0{,}2483
\end{equation}
Sau khi phân bổ $30$ câu ngoài phạm vi, trên $210$ câu hỏi có nhãn nguồn, tỷ lệ này xác định chính xác $52$ câu hỏi hình ảnh và $158$ câu hỏi văn bản (Bảng~\ref{tab:testset_dist}).

\begin{algorithm2e}[H]
\caption{Quy trình khởi tạo bộ dữ liệu kiểm thử ngẫu nhiên phân tầng}
\label{alg:gen}
\SetAlgoLined
\KwIn{Tổng quy mô $N = 240$; số câu ngoài phạm vi $N_{op} = 30$; hạt giống ngẫu nhiên \texttt{seed}}
\KwOut{Tập dữ liệu \texttt{(câu\_hỏi, đáp\_án\_chuẩn, phân\_loại, sách\_nguồn, trang\_nguồn, nhãn\_hình)}}
Xác định số lượng đoạn văn bản hợp lệ $n_{chunk}$ và số khung hình $n_{anh}$ trong chỉ mục\;
$p_{hinh} \leftarrow n_{anh} / (n_{chunk} + n_{anh})$\;
$N' \leftarrow N - N_{op}$; $\quad N_{hinh} \leftarrow \mathrm{round}(N' \cdot p_{hinh})$; $\quad N_{vb} \leftarrow N' - N_{hinh}$\;
$rows \leftarrow \emptyset$\;
\ForEach{$(loai, N_{loai}, pool) \in \{(\text{văn bản}, N_{vb}, \mathcal{P}_{vb}),\ (\text{hình}, N_{hinh}, \mathcal{P}_{anh})\}$}{
    Lấy mẫu ngẫu nhiên không hoàn lại $N_{loai}$ phần tử từ $pool$\;
    \ForEach{phần tử dữ liệu $x$}{
        $(q, gt) \leftarrow \text{Khởi\_tạo\_câu\_hỏi\_từ\_ngữ\_cảnh}(x)$\;
        Thêm \texttt{(q, gt, loai, source(x), page(x))} vào danh sách $rows$\;
    }
}
\ForEach{$i \in 1..N_{op}$}{
    $mon \leftarrow$ chọn ngẫu nhiên một môn học ngoài KHTN (Toán, Văn, Sử, Tin)\;
    $(q, gt) \leftarrow \text{Khởi\_tạo\_câu\_hỏi\_ngoại\_vi}(mon)$ \tcp*{$gt$: chỉ dẫn từ chối}
    Thêm \texttt{(q, gt, ngoai\_pham\_vi, -, -)} vào danh sách $rows$\;
}
\Return $rows$ (chuyển qua bước chuyên gia rà soát và kiểm duyệt xác nhận)\;
\end{algorithm2e}

Quy trình tích hợp cơ chế ngắt mạch an toàn, tự động dừng tác vụ nếu tỷ lệ lỗi gọi mô hình vượt quá ngưỡng $30\%$ trong cửa sổ trượt $20$ yêu cầu liên tiếp. Toàn bộ $240$ câu hỏi sau khi khởi tạo đều trải qua bước chuyên gia rà soát thủ công toàn diện trên văn bản để hiệu chỉnh ngữ nghĩa và chuẩn hoá nhãn trước khi cấp cờ xác thực \texttt{human\_reviewed: true} trong tệp cấu hình \texttt{meta.json}. Thời gian rà soát trung bình đạt $\approx 48{,}7$ giây/câu hỏi.

\begin{table}[H]
\centering
\caption{Phân bố bộ dữ liệu kiểm thử 240 câu theo bộ sách giáo khoa và nhóm câu hỏi (Nguồn: Bộ dữ liệu kiểm thử của đồ án)}
\label{tab:testset_dist}
\small
\begin{adjustbox}{max width=\textwidth}
\begin{tabular}{lccc}
\toprule
\textbf{Bộ sách giáo khoa} & \textbf{Câu hỏi văn bản} & \textbf{Câu hỏi hình ảnh} & \textbf{Tổng số câu} \\
\midrule
Cánh Diều (4 cuốn)            & $52$ & $16$ & $68$ \\
Chân trời sáng tạo (4 cuốn)   & $52$ & $20$ & $72$ \\
Kết nối tri thức (4 cuốn)     & $54$ & $16$ & $70$ \\
\midrule
\textbf{Tổng số câu có nhãn nguồn} & \textbf{158} & \textbf{52} & \textbf{210} \\
Nhóm câu hỏi ngoài phạm vi (không nhãn nguồn) & \multicolumn{2}{c}{-} & $30$ \\
\midrule
\textbf{Tổng quy mô bộ kiểm thử} & & & \textbf{240} \\
\bottomrule
\end{tabular}
\end{adjustbox}
\end{table}

\begin{figure}[H]
    \centering
    \includegraphics[width=0.88\textwidth]{src/images/chapter4/phan_bo_loai_cau_hoi.png}
    \caption{Thành phần phân bổ bộ dữ liệu kiểm thử 240 câu hỏi theo từng nhóm loại (Nguồn: Bộ dữ liệu kiểm thử của đồ án)}
    \label{fig:phan_bo_loai}
\end{figure}

Hình~\ref{fig:phan_bo_loai} trực quan hóa tỷ trọng giữa ba nhóm câu hỏi: câu hỏi văn bản chiếm đa số ($158$ câu, tương đương $65{,}8\%$), phản ánh khối lượng tri thức chủ đạo của bộ sách giáo khoa; câu hỏi hình ảnh chiếm $52$ câu ($21{,}7\%$); và nhóm câu hỏi ngoài phạm vi gồm $30$ câu ($12{,}5\%$) giữ vai trò kiểm thử năng lực từ chối hiện tượng ảo giác.

Bảng~\ref{tab:sample_q} minh họa một số câu hỏi tiêu biểu thuộc ba nhóm dữ liệu kiểm thử cùng đáp án chuẩn và nhãn nguồn trích xuất tương ứng.

\begin{table}[H]
\centering
\caption{Một số mẫu câu hỏi kiểm thử tiêu biểu kèm nhãn nguồn xác thực (Nguồn: Trích xuất từ bộ dữ liệu kiểm thử)}
\label{tab:sample_q}
\small
\begin{adjustbox}{max width=\textwidth}
\begin{tabular}{p{0.36\textwidth}p{0.44\textwidth}p{0.20\textwidth}}
\toprule
\textbf{Nội dung câu hỏi} & \textbf{Đáp án chuẩn đối sánh} & \textbf{Nguồn trích dẫn} \\
\midrule
Theo đoạn văn, quá trình quang hợp ở thực vật đã chuyển hóa những chất nào để tạo ra glucose? & Khí carbon dioxide và nước. & Lớp 9 Kết nối tri thức, trang 135 \\
\addlinespace
Theo đoạn văn, để ánh sáng truyền trong sợi quang gặp lớp vỏ và bị phản xạ toàn phần, chiết suất của lõi và lớp vỏ phải có đặc điểm gì? & Chiết suất của lõi phải lớn hơn chiết suất của lớp vỏ. & Lớp 9 Cánh Diều, trang 23 \\
\addlinespace
Dựa vào Hình 24.3, hãy liệt kê ba vai trò của đa dạng sinh học đối với con người được minh họa trong hình. & Làm nơi tham quan, học tập; Cung cấp nguồn lương thực thực phẩm; Làm dược liệu. & Lớp 6 Cánh Diều, trang 132 (Hình 24.3) \\
\addlinespace
Trong lập trình Python, cú pháp nào được sử dụng để in ra màn hình chuỗi ``Hello World''? & \textit{Câu hỏi ngoài phạm vi môn học:} Hệ thống thông báo không tìm thấy thông tin trong sách giáo khoa và từ chối trả lời bằng tri thức ngoài học liệu chuẩn. & Ngoài phạm vi \\
\bottomrule
\end{tabular}
\end{adjustbox}
\end{table}

\section{Phương pháp luận và hệ thống độ đo đánh giá}
Do đầu ra của hệ thống trợ lý ảo là ngôn ngữ tự nhiên mở, việc chỉ áp dụng các thang đo so khớp từ khóa hình thức truyền thống như BLEU~\cite{papineni2002bleu} hay ROUGE~\cite{lin2004rouge} là không đầy đủ, bởi một diễn đạt khoa học chính xác có thể được trình bày bằng nhiều cấu trúc ngữ pháp tương đương. Đề tài xây dựng khung đánh giá hai pha toàn diện:
\begin{enumerate}
    \item \textbf{Pha 1: Đánh giá năng lực truy xuất.} Đo lường khách quan bằng các chỉ số truy xuất thông tin chuẩn mực dựa trên sự đối sánh chính xác với nhãn nguồn;
    \item \textbf{Pha 2: Đánh giá chất lượng sinh ngôn ngữ.} Ứng dụng phương pháp đánh giá bằng mô hình ngôn ngữ lớn làm giám khảo độc lập~\cite{zheng2023llmjudge, liu2023geval, es2024ragas, min2023factscore, saadfalcon2023ares} với mô hình ngôn ngữ độc lập chấm điểm định lượng theo ba tiêu chí cốt lõi.
\end{enumerate}

Toàn bộ quy trình đánh giá tự động hai pha được khái quát trong Thuật toán~\ref{alg:eval} và minh họa tại Hình~\ref{fig:eval_pipeline_concept}.

\subsection{Độ đo đánh giá pha truy xuất tri thức}
\label{sec:retrieval_eval_def}
Xét câu hỏi kiểm thử với danh sách $k$ phân đoạn văn bản do hệ thống truy xuất được $R_k = (d_1, d_2, \dots, d_m)$ ($m \leq k$) được sắp xếp giảm dần theo điểm số liên quan. Gọi $G$ là tập hợp toàn bộ các phân đoạn trong chỉ mục thuộc đúng trang sách chuẩn đối sánh $(B^{*}, p^{*})$ của câu hỏi đó, với quy mô $m^{*} = |G|$.

Khác với chỉ số độ trúng nhị phân Hit@$k$, đề tài chuẩn hóa các độ đo theo tỷ lệ phân đoạn thực tế của trang nguồn:
\begin{align}
\mathrm{Precision@}k &= \frac{|R_k \cap G|}{k}, \quad \mathrm{Recall@}k = \frac{|R_k \cap G|}{m^{*}}, \\
\mathrm{F1@}k &= \frac{2 \cdot \mathrm{Precision@}k \cdot \mathrm{Recall@}k}{\mathrm{Precision@}k + \mathrm{Recall@}k}.
\end{align}

Chỉ số \textbf{F1@$k$ được tính toán theo phương pháp Macro}: giá trị F1@$k$ được xác định độc lập cho từng câu hỏi trước khi tính trung bình cộng toàn bộ tập kiểm thử. Cách tiếp cận này đảm bảo phản ánh trung thực các trường hợp truy xuất thất bại ($Recall = 0$), tuân thủ chặt chẽ bất đẳng thức Jensen.

Chỉ số nghịch đảo thứ hạng trung bình MRR~\cite{voorhees1999trec} đo lường chất lượng vị trí của phân đoạn đúng đầu tiên:
\begin{equation}
\mathrm{MRR} = \frac{1}{N}\sum_{i=1}^{N} \frac{1}{\mathrm{rank}_i}
\end{equation}
với $\mathrm{rank}_i = \infty$ (đóng góp giá trị $0$) nếu không có phân đoạn chuẩn nào xuất hiện trong danh sách $R_k$.

Cần nhấn mạnh rằng chỉ số Precision@$k$ bị chặn trên bởi giới hạn trần lý thuyết $\min(k, m^{*})/k$. Trong trường hợp trang sách chuẩn chỉ chứa $m^{*} = 3$ đoạn văn bản, giá trị Precision@5 tối đa chỉ có thể đạt $0{,}6$ ngay cả khi hệ thống truy xuất hoàn hảo. Do đó, đề tài báo cáo kèm giá trị trần lý thuyết (\texttt{tranP@k}) để có góc nhìn đối sánh chuẩn xác. Điều kiện so khớp chuẩn luôn là \textbf{chính xác tuyệt đối về số trang in}.

\subsection{Độ đo đánh giá chất lượng sinh văn bản bằng mô hình giám khảo}
\label{sec:judge_pha2}
Câu trả lời do mô hình Qwen2.5 sinh ra được đánh giá độc lập bởi mô hình ngôn ngữ lớn đóng vai trò giám khảo thông qua nền tảng Groq API. Mô hình giám khảo tiếp nhận ngữ cảnh truy xuất, câu hỏi, câu trả lời sinh ra và đáp án chuẩn để chấm điểm trên thang đo từ $1$ đến $5$ theo ba tiêu chí chuẩn hóa:
\begin{itemize}
    \item \textbf{Tính đúng:} Đánh giá mức độ chính xác về mặt nội dung khoa học của câu trả lời khi so sánh với đáp án chuẩn.
    \item \textbf{Độ trung thực:} Đánh giá mức độ bám sát của câu trả lời đối với ngữ cảnh tài liệu được cung cấp, kiểm soát hiện tượng tự suy diễn hoặc sinh ảo giác thông tin.
    \item \textbf{Độ liên quan:} Đánh giá mức độ tập trung, trả lời trực diện vào trọng tâm câu hỏi của người học, không lan man.
\end{itemize}

Quá trình đánh giá sử dụng tập hợp bốn mô hình giám khảo luân phiên gồm:
\begin{itemize}
    \item \texttt{qwen/qwen3.8-27b};
    \item \texttt{qwen/qwen3.6-27b};
    \item \texttt{openai/gpt-oss-120b};
    \item \texttt{openai/gpt-oss-20b}.
\end{itemize}
Quy trình gửi yêu cầu chấm điểm tích hợp các phương pháp thử lại (retry) gồm:
\begin{itemize}
    \item Thử lại tối đa ba lần đối với các lỗi tạm thời (mã trạng thái 429 hoặc lỗi định dạng dữ liệu phản hồi);
    \item Áp dụng khoảng thời gian chờ lũy tiến giữa các lượt thử lại liên tiếp;
    \item Tự động luân chuyển sang mô hình dự phòng kế tiếp trong danh sách khi yêu cầu thử lại vượt quá số lần quy định, đảm bảo $100\%$ các câu hỏi kiểm thử đều được hoàn tất chấm điểm.
\end{itemize}

\subsection{Trục phân tích tổng hợp theo loại câu hỏi}
\label{sec:truc_tong_hop}
Kết quả thực nghiệm được tổng hợp và phân tích chuyên sâu theo \textbf{loại câu hỏi} (văn bản, hình ảnh, ngoài phạm vi) thay vì phân tách riêng rẽ theo từng cuốn sách. Phương pháp này đảm bảo tính vững chắc về mặt thống kê khi áp dụng lấy mẫu ngẫu nhiên không hoàn lại, phản ánh chính xác năng lực cốt lõi của hệ thống trước các dạng thức tương tác khác nhau của học sinh.

\begin{algorithm2e}[H]
\caption{Khung quy trình thực nghiệm đánh giá tự động hai pha}
\label{alg:eval}
\SetAlgoLined
\KwIn{Bộ kiểm thử $240$ câu hỏi; Mô hình giám khảo độc lập $J$}
\KwOut{Hệ thống chỉ số truy xuất thông tin; Điểm số đánh giá pha sinh ngôn ngữ}
\ForEach{câu hỏi $q$ (thuộc phân loại $\ell$, nhãn nguồn $(B^{*}, p^{*})$)}{
    $R \leftarrow \text{HybridRetriever}(q)$ \tcp*{Truy xuất ứng viên theo luồng thực tế}
    \If{$\ell \neq$ ngoài phạm vi}{
        Tính toán các chỉ số Precision@k, Recall@k, F1@k, MRR dựa trên $R$ và $(B^{*}, p^{*})$\;
    }
    $ctx \leftarrow$ tổng hợp ngữ cảnh từ các phân đoạn trong $R$\;
    $ans \leftarrow \text{Qwen2.5}(\text{prompt}(ctx, q))$\;
    $(corr, faith, rel) \leftarrow J(q, gt, ctx, ans)$ \tcp*{Mô hình giám khảo độc lập chấm điểm}
    Ghi nhận toàn bộ chỉ số thành phần gắn với định danh câu hỏi\;
}
\Return Tổng hợp chỉ số IR theo cấu hình truy xuất và điểm số giám khảo theo nhóm câu hỏi $\ell$\;
\end{algorithm2e}

\begin{figure}[H]
    \centering
    \includegraphics[width=0.96\textwidth]{src/images/chapter4/quy_trinh_danh_gia_hai_pha.png}
    \caption{Sơ đồ quy trình thực nghiệm đánh giá tự động hai pha độc lập (Nguồn: Đề xuất của đồ án)}
    \label{fig:eval_pipeline_concept}
\end{figure}

\section{Kết quả thực nghiệm định lượng}
\label{sec:ketqua}
Toàn bộ các chỉ số dưới đây được tính toán dựa trên phương pháp trung bình cộng độc lập trên từng câu hỏi, đảm bảo tính đại diện khách quan cho toàn bộ tập dữ liệu kiểm thử.

\subsection{Thực nghiệm đối chiếu các cấu hình truy xuất}
\label{sec:ablation}
Bằng việc thiết lập ba công tắc độc lập (kênh truy xuất, cơ chế tái xếp hạng, cổng lọc khoảng cách tương đối), hệ thống tiến hành đánh giá thực nghiệm bóc tách thành phần trên $210$ câu hỏi có nhãn nguồn ở quy mô vận hành thực tế ($20$ ứng viên ban đầu cho mỗi kênh). Bảng~\ref{tab:ablation_main} và Hình~\ref{fig:so_sanh_truy_xuat} so sánh trực tiếp bốn phương pháp truy xuất cốt lõi theo mục tiêu thiết kế và bài toán nghiên cứu của đề tài.

\begin{table}[H]
\centering
\caption{Kết quả đối chiếu bốn phương pháp truy xuất ở quy mô vận hành thực tế ($210$ câu hỏi có nhãn nguồn) (Nguồn: Kết quả thực nghiệm của đồ án)}
\label{tab:ablation_main}
\small
\begin{adjustbox}{max width=\textwidth}
\begin{tabular}{lccccc}
\toprule
\textbf{Phương pháp truy xuất} & \textbf{MRR} & \textbf{R@1} & \textbf{R@3} & \textbf{R@5} & \textbf{R@10} \\
\midrule
Chỉ mục thưa BM25 thuần & 0,6636 & 0,0905 & 0,1475 & 0,1796 & 0,2303 \\
Truy xuất ngữ nghĩa véc-tơ đặc & 0,5664 & 0,0647 & 0,1338 & 0,1778 & 0,2302 \\
Truy xuất lai (không tái xếp hạng) & 0,6947 & 0,0914 & 0,1609 & 0,2031 & \textbf{0,2643} \\
\textbf{Truy xuất lai kết hợp tái xếp hạng (Đề xuất)} & \textbf{0,7789} & \textbf{0,1041} & \textbf{0,1804} & \textbf{0,2085} & 0,2458 \\
\bottomrule
\end{tabular}
\end{adjustbox}
\end{table}

\begin{figure}[H]
    \centering
    \includegraphics[width=0.92\textwidth]{src/images/chapter4/so_sanh_truy_xuat.png}
    \caption{So sánh định lượng hiệu năng truy xuất giữa bốn phương pháp theo các chỉ số MRR, R@1, R@3, R@5 và R@10 (Nguồn: Kết quả thực nghiệm của đồ án)}
    \label{fig:so_sanh_truy_xuat}
\end{figure}

Từ kết quả định lượng tại Bảng~\ref{tab:ablation_main}, hai phát hiện khoa học quan trọng được rút ra:

\textbf{(i) Hiệu năng vượt trội của cấu hình đề xuất và hiện tượng suy giảm nhẹ tại R@10:} Cấu hình truy xuất lai kết hợp tái xếp hạng bằng mô hình tương tác chéo đạt hiệu năng cao nhất trên hầu hết các chỉ số chủ đạo: chỉ số MRR tăng mạnh từ $0{,}6947$ lên $\mathbf{0{,}7789}$ (vượt trội so với mức $0{,}6636$ của chỉ mục thưa BM25 thuần và $0{,}5664$ của kênh ngữ nghĩa thuần), đồng thời chỉ số R@1 cải thiện từ $0{,}0914$ lên $\mathbf{0{,}1041}$. Tuy nhiên, ở chỉ số R@10, cấu hình đề xuất lại ghi nhận mức giảm nhẹ ($0{,}2458$ so với $0{,}2643$). Giả thuyết khoa học hợp lý cho hiện tượng này là việc áp dụng ngưỡng điểm số tối thiểu của bộ tái xếp hạng (\texttt{RERANK\_SCORE\_MIN = 0.59}) đã vô tình loại bỏ một số phân đoạn tài liệu đúng nhưng bị đẩy xuống các vị trí cận biên ngoài top-10.

\textbf{(ii) Tác động bất lợi của cổng lọc khoảng cách tương đối ở quy mô thực tế:} Khi kích hoạt cổng lọc theo khoảng cách tương đối, chỉ số MRR của cấu hình đề xuất suy giảm rõ rệt từ $0{,}7779$ xuống $0{,}7157$, và R@10 giảm từ $0{,}2410$ xuống $0{,}1974$. Nguyên nhân bắt nguồn từ giả định cứng nhắc của cổng lọc khi cho rằng tài liệu đúng luôn phải nằm rất gần tài liệu có điểm số cao nhất, điều này không phù hợp với các câu hỏi phức hợp đòi hỏi tổng hợp kiến thức từ nhiều đoạn văn bản khác nhau. Do đó, việc vô hiệu hóa cổng lọc khoảng cách tương đối là quyết định tối ưu đã được kiểm chứng bằng thực nghiệm.

Bảng~\ref{tab:ablation_rerank} phân tích chi tiết đóng góp độc lập của bước tái xếp hạng bằng mô hình tương tác chéo trên từng kênh truy xuất riêng biệt (ở quy mô thử nghiệm $50$ ứng viên ban đầu, tắt cổng lọc khoảng cách).

\begin{table}[H]
\centering
\caption{Đóng góp của bước tái xếp hạng bằng mô hình tương tác chéo trên các kênh truy xuất ($210$ câu hỏi) (Nguồn: Kết quả thực nghiệm của đồ án)}
\label{tab:ablation_rerank}
\small
\begin{adjustbox}{max width=\textwidth}
\begin{tabular}{lccc}
\toprule
\textbf{Kênh truy xuất thành phần} & \textbf{Trạng thái tái xếp hạng} & \textbf{Chỉ số MRR} & \textbf{Chỉ số R@1} \\
\midrule
\multirow{2}{*}{Truy xuất lai} & Tắt & 0,6947 & 0,0914 \\
 & \textbf{Kích hoạt} & \textbf{0,7789} & \textbf{0,1041} \\
\midrule
\multirow{2}{*}{Chỉ mục từ khóa BM25} & Tắt & 0,6636 & 0,0905 \\
 & Kích hoạt & 0,7676 & 0,1026 \\
\midrule
\multirow{2}{*}{Kênh ngữ nghĩa véc-tơ đặc} & Tắt & 0,5664 & 0,0647 \\
 & Kích hoạt & 0,7457 & 0,0994 \\
\bottomrule
\end{tabular}
\end{adjustbox}
\end{table}

Bước tái xếp hạng giúp nâng cao chỉ số MRR trên cả ba kênh truy xuất. Đáng chú ý, kênh truy xuất ngữ nghĩa theo véc-tơ đặc, vốn có hiệu năng cơ sở thấp nhất khi đứng độc lập, lại ghi nhận mức tăng trưởng tương đối mạnh mẽ nhất ở chỉ số R@1 (tăng từ $0{,}0647$ lên $0{,}0994$, tương ứng mức cải thiện $\mathbf{+53{,}6\%}$, so với $+13{,}4\%$ của BM25 và $+13{,}9\%$ của kênh lai). Kết quả này khẳng định: đối với tập học liệu khoa học chứa mật độ thuật ngữ và ký hiệu cao, mô hình nhúng hai nhánh bắt buộc phải được kết hợp cùng bước tái xếp hạng bằng mô hình tương tác chéo có khả năng phân tích tương tác từ ngữ sâu sắc.

\subsection{Tổng hợp các chỉ số hiệu năng hệ thống}
Bảng~\ref{tab:summary} tổng hợp toàn diện các chỉ số thực nghiệm của cấu hình đề xuất tối ưu (truy xuất lai kết hợp tái xếp hạng bằng mô hình tương tác chéo, tắt cổng lọc khoảng cách tương đối) ở quy mô vận hành thực tế.

\begin{table}[H]
\centering
\caption{Tổng hợp các chỉ số hiệu năng của cấu hình hệ thống đề xuất (Nguồn: Kết quả thực nghiệm của đồ án)}
\label{tab:summary}
\small
\begin{adjustbox}{max width=\textwidth}
\begin{tabular}{lc}
\toprule
\textbf{Chỉ số đánh giá} & \textbf{Giá trị định lượng} \\
\midrule
\multicolumn{2}{l}{\textit{Pha 1: Hiệu năng truy xuất (210 câu hỏi có nhãn nguồn)}} \\
Chỉ số MRR & \textbf{0,7789} \\
Độ phủ Recall@1 & 0,1041 \\
Độ phủ Recall@3 & 0,1804 \\
Độ phủ Recall@5 & 0,2085 \\
Độ phủ Recall@10 & 0,2458 \\
Độ chính xác Precision@5 (Trần lý thuyết \texttt{tranP@5} = 0,9867) & 0,2810 \\
\midrule
\multicolumn{2}{l}{\textit{Khả năng kiểm soát phạm vi (30 câu hỏi ngoài phạm vi)}} \\
Tỷ lệ nhận diện và từ chối trả lời chính xác & \textbf{0,9667} ($29/30$ câu) \\
\midrule
\multicolumn{2}{l}{\textit{Pha 2: Chất lượng sinh ngôn ngữ (Mô hình giám khảo độc lập, thang điểm 1 - 5)}} \\
Tính đúng & 4,033 / 5 \\
Độ trung thực & 4,417 / 5 \\
Độ liên quan & 4,492 / 5 \\
\bottomrule
\end{tabular}
\end{adjustbox}
\end{table}

\subsection{Chất lượng câu trả lời sinh ra theo từng nhóm câu hỏi}
\label{sec:judge_loai}
Bảng~\ref{tab:judge_loai} và Biểu đồ~\ref{fig:judge} trình bày chi tiết điểm số trung bình do hội đồng mô hình giám khảo độc lập đánh giá trên $240$ câu hỏi kiểm thử, phân tách theo ba nhóm câu hỏi chính.

\begin{table}[H]
\centering
\caption{Chất lượng câu trả lời theo từng nhóm câu hỏi do mô hình giám khảo độc lập chấm điểm (Nguồn: Kết quả thực nghiệm của đồ án)}
\label{tab:judge_loai}
\small
\begin{adjustbox}{max width=\textwidth}
\begin{tabular}{lcccc}
\toprule
\textbf{Nhóm câu hỏi kiểm thử} & \textbf{Số lượng câu} & \textbf{Tính đúng / 5} & \textbf{Độ trung thực / 5} & \textbf{Độ liên quan / 5} \\
\midrule
Câu hỏi văn bản & 158 & 4,19 & 4,56 & 4,59 \\
Câu hỏi hình ảnh & 52 & 3,08 & 3,65 & 3,90 \\
Câu hỏi ngoài phạm vi & 30 & 4,87 & 5,00 & 5,00 \\
\bottomrule
\end{tabular}
\end{adjustbox}
\end{table}

\begin{figure}[H]
    \centering
    \includegraphics[width=0.92\textwidth]{src/images/chapter4/judge_scores_theo_loai.png}
    \caption{Phân bố điểm số chất lượng câu trả lời của Qwen2.5 theo từng nhóm câu hỏi (Nguồn: Kết quả thực nghiệm của đồ án)}
    \label{fig:judge}
\end{figure}

Phân tích định lượng từ kết quả giám khảo cho thấy ba đặc trưng nổi bật:
\begin{itemize}
    \item \textbf{Khả năng nhận diện câu hỏi ngoài phạm vi đạt mức gần như tuyệt đối:} Điểm số nhóm câu hỏi ngoài phạm vi đạt $4{,}87/5$ về Tính đúng, cùng điểm tuyệt đối $5{,}00/5$ về Độ trung thực và Độ liên quan. Đây là minh chứng định lượng rõ nét cho thấy hệ thống tuân thủ chặt chẽ nguyên tắc bảo toàn tri thức: khi câu hỏi không thuộc phạm vi học liệu Khoa học tự nhiên, hệ thống từ chối trả lời chính xác trong $96{,}67\%$ trường hợp ($29/30$ câu), triệt tiêu nguy cơ tự suy đoán bằng tri thức ngoài sách giáo khoa.
    \item \textbf{Nhóm câu hỏi hình ảnh đạt kết quả khiêm tốn hơn:} Điểm Tính đúng của nhóm câu hỏi hình ảnh ($3{,}08/5$) thấp hơn rõ rệt so với nhóm văn bản ($4{,}19/5$), phản ánh những thách thức cố hữu trong việc liên kết biểu diễn thị giác và văn bản trong sách giáo khoa.
    \item \textbf{Mối tương quan giữa Tính đúng và Độ trung thực:} Trong cả ba nhóm, điểm Tính đúng luôn thấp hơn điểm Độ trung thực ($4{,}19 < 4{,}56$; $3{,}08 < 3{,}65$). Điều này phản ánh xu hướng nhất quán của mô hình ngôn ngữ: khi ngữ cảnh trích xuất bị thiếu hụt hoặc chưa hoàn hảo, mô hình vẫn duy trì tính trung thực cao với ngữ cảnh nhận được hơn là tự ý suy diễn để hoàn thiện đáp án.
\end{itemize}

Bảng~\ref{tab:case_study} minh họa hai trường hợp thực tế tiêu biểu trích xuất từ nhật ký kiểm thử: một ca trả lời câu hỏi học liệu Khoa học tự nhiên kèm số trang trích dẫn chính xác và một ca nhận diện câu hỏi ngoài phạm vi để từ chối trả lời, chứng minh năng lực kiểm soát ảo giác của hệ thống.

\begin{table}[H]
\centering
\caption{Minh họa hai ca kiểm thử thực tế về phản hồi của hệ thống và đánh giá của giám khảo độc lập (Nguồn: Trích xuất từ nhật ký kiểm thử của đồ án)}
\label{tab:case_study}
\small
\begin{adjustbox}{max width=\textwidth}
\begin{tabular}{p{3.5cm}p{11.5cm}}
\toprule
\textbf{Thành phần đối sánh} & \textbf{Nội dung thực tế ghi nhận} \\
\midrule
\multicolumn{2}{l}{\textbf{Trường hợp 1: Câu hỏi học liệu Khoa học tự nhiên (Kiểm thử năng lực trích dẫn và trả lời có căn cứ)}} \\
\addlinespace[2pt]
Câu hỏi người học & \textit{``Tại sao lá cây lại có màu xanh lục và bào quan nào thực hiện chức năng quang hợp?''} \\
Ngữ cảnh trích xuất từ SGK & Phân đoạn từ sách \textbf{KHTN 6 Cánh Diều, trang 96}: \textit{``Lục lạp chứa chất diệp lục có khả năng hấp thụ năng lượng ánh sáng mặt trời để tổng hợp chất hữu cơ... Chất diệp lục phản xạ ánh sáng màu lục nên mắt ta nhìn thấy lá có màu xanh.''} \\
Phản hồi của Qwen2.5 & \textit{``Lá cây có màu xanh lục vì trong tế bào lá có chứa bào quan lục lạp. Bên trong lục lạp chứa chất diệp lục; sắc tố này hấp thụ ánh sáng mặt trời và phản xạ ánh sáng màu lục vào mắt ta. Lục lạp chính là bào quan thực hiện chức năng quang hợp của cây xanh. (Trích dẫn: KHTN 6 Cánh Diều, trang 96).''} \\
Chấm điểm của Giám khảo & \textbf{Tính đúng: 5/5} $\mid$ \textbf{Độ trung thực: 5/5} $\mid$ \textbf{Độ liên quan: 5/5} \\
Nhận xét của Giám khảo & \textit{Câu trả lời chính xác, bám sát tuyệt đối ngữ cảnh tài liệu, trích dẫn đúng số trang và không suy diễn thông tin ngoài sách giáo khoa.} \\
\midrule
\multicolumn{2}{l}{\textbf{Trường hợp 2: Câu hỏi ngoài phạm vi môn học (Kiểm thử năng lực phòng chống ảo giác kiến thức)}} \\
\addlinespace[2pt]
Câu hỏi người học & \textit{``Chiến dịch Điện Biên Phủ diễn ra vào năm nào và có ý nghĩa lịch sử gì đối với dân tộc Việt Nam?''} \\
Ngữ cảnh trích xuất từ SGK & \textit{Không tìm thấy phân đoạn tài liệu phù hợp trong cơ sở tri thức sách giáo khoa KHTN (điểm tương đồng dưới ngưỡng).} \\
Phản hồi của Qwen2.5 & \textit{``Xin lỗi, câu hỏi này thuộc lĩnh vực Lịch sử, nằm ngoài phạm vi cơ sở tri thức sách giáo khoa Khoa học tự nhiên (Lớp 6 đến Lớp 9) mà hệ thống được cung cấp. Tôi chỉ có thể hỗ trợ giải đáp các kiến thức thuộc môn học Khoa học tự nhiên.''} \\
Chấm điểm của Giám khảo & \textbf{Tính đúng: 5/5} $\mid$ \textbf{Độ trung thực: 5/5} $\mid$ \textbf{Độ liên quan: 5/5} \\
Nhận xét của Giám khảo & \textit{Hệ thống nhận diện chuẩn xác câu hỏi ngoài phạm vi, từ chối trả lời đúng quy định sư phạm và không tự ý bịa đặt câu trả lời.} \\
\bottomrule
\end{tabular}
\end{adjustbox}
\end{table}

\section{Thảo luận chuyên sâu về các phát hiện khoa học}
\label{sec:thaoluan}
Tổng hợp từ các kết quả thực nghiệm định lượng, đề tài rút ra năm nhận định khoa học then chốt:

\begin{enumerate}
    \item \textbf{Cơ chế tái xếp hạng bằng mô hình tương tác chéo là thành phần quyết định hiệu năng truy xuất:} Việc kết hợp truy xuất lai và mô hình tương tác chéo mang lại sự vượt trội về chỉ số MRR ($0{,}7789$) và các độ phủ ở thứ hạng cao ($R@1$, $R@3$, $R@5$). Đồng thời, thực nghiệm chứng minh cổng lọc theo khoảng cách tương đối gây tác động tiêu cực đến độ phủ ở quy mô dữ liệu thực tế, khẳng định việc sử dụng ngưỡng điểm tuyệt đối của bộ tái xếp hạng là giải pháp tối ưu.
    \item \textbf{Kênh ngữ nghĩa véc-tơ đặc chịu sự suy thoái trước ký hiệu khoa học và cần sự hỗ trợ của mô hình tái xếp hạng:} Đối với miền dữ liệu chứa mật độ cao công thức và thuật ngữ chuyên ngành, kiến trúc hai nhánh mã hóa độc lập (kích thước không gian nhúng $1024$ chiều) gặp khó khăn trong việc phân biệt các ký tự định danh chính xác, dẫn tới chỉ số R@1 ban đầu thấp ($0{,}0647$). Bước tái xếp hạng đã phát huy tối đa hiệu quả khi nâng chỉ số này lên $0{,}0994$ (tăng $+53{,}6\%$).
    \item \textbf{Năng lực kiểm soát ảo giác được chứng minh thông qua nhóm câu hỏi ngoài phạm vi:} Tỷ lệ từ chối trả lời chính xác $0{,}9667$ cùng điểm số giám khảo gần tuyệt đối trên các câu hỏi ngoài phạm vi khẳng định cấu trúc prompt sư phạm nghiêm ngặt và cơ chế RAG của hệ thống đã kiểm soát hiệu quả hiện tượng mô hình tự phát sinh tri thức ngoài học liệu chuẩn.
    \item \textbf{Nút thắt kỹ thuật ở nhóm câu hỏi hình ảnh tập trung tại khâu liên kết ngữ cảnh:} Chất lượng câu trả lời của nhóm hình ảnh ($3{,}08/5$) thấp hơn nhóm văn bản, đòi hỏi các nghiên cứu tiếp theo cần tập trung nâng cao năng lực biểu diễn đa phương thức ở cả hai phía truy xuất và sinh văn bản.
    \item \textbf{Đánh đổi giữa độ chính xác công thức và tính toàn vẹn ngôn ngữ trong OCR:} Thực nghiệm đối sánh độc lập khẳng định Tesseract nhận dạng tốt văn bản tiếng Việt ($93\%$ nội dung) nhưng làm hỏng chỉ số dưới của công thức, trong khi mô hình thị giác - ngôn ngữ đọc tốt công thức nhưng lại làm tăng tỷ lệ lỗi dấu thanh tiếng Việt. Giải pháp kiến trúc lai phân cấp theo dòng đã giải quyết thành công sự đánh đổi này trên toàn bộ cơ sở tri thức ($4\,293$ lượt ghép dòng thành công, $0$ lỗi ngoại lệ).
\end{enumerate}

\subsection{Đối chiếu thực nghiệm giữa RAG đa phương thức và RAG thuần văn bản}
\label{sec:mm}
Nhằm thực hiện đầy đủ mục tiêu nghiên cứu đối chiếu, hệ thống tiến hành thực nghiệm so sánh ghép cặp trực tiếp trên đúng $52$ câu hỏi hình ảnh có nhãn chuẩn sinh từ siêu dữ liệu hình ảnh xác thực. Bảng~\ref{tab:m2c} và Hình~\ref{fig:so_sanh_m2c} đối chiếu chất lượng câu trả lời giữa cấu hình chỉ sử dụng ngữ cảnh văn bản và cấu hình có bổ sung ngữ cảnh đa phương thức gồm chuỗi siêu dữ liệu hình ảnh trích xuất từ trang sách.

\begin{table}[H]
\centering
\caption{Kết quả đối chiếu thực nghiệm đa phương thức trên 52 câu hỏi hình ảnh (Nguồn: Kết quả thực nghiệm của đồ án)}
\label{tab:m2c}
\small
\begin{adjustbox}{max width=\textwidth}
\begin{tabular}{lccc}
\toprule
\textbf{Tiêu chí đánh giá} & \textbf{Cấu hình thuần văn bản} & \textbf{Cấu hình đa phương thức} & \textbf{Mức chênh lệch ($\Delta$)} \\
\midrule
Tính đúng & 3,077 & 3,000 & $-0{,}077$ \\
Độ trung thực & 3,654 & 3,615 & $-0{,}038$ \\
Độ liên quan & 3,904 & 3,865 & $-0{,}038$ \\
\bottomrule
\end{tabular}
\end{adjustbox}
\end{table}

\begin{figure}[H]
    \centering
    \includegraphics[width=0.82\textwidth]{src/images/chapter4/so_sanh_da_phuong_thuc.png}
    \caption{Trực quan hóa mức chênh lệch chất lượng câu trả lời giữa cấu hình thuần văn bản và đa phương thức trên 52 câu hỏi hình ảnh (Nguồn: Kết quả thực nghiệm của đồ án)}
    \label{fig:so_sanh_m2c}
\end{figure}

Kết quả đối chiếu từng câu hỏi ghi nhận: $46/52$ câu ($88{,}5\%$) giữ nguyên điểm số, $5$ câu ghi nhận mức giảm $1$ điểm và $1$ câu ghi nhận mức tăng $1$ điểm. Cả ba tiêu chí chất lượng đều có sự suy giảm nhẹ khi bổ sung ngữ cảnh đa phương thức dạng chuỗi mô tả vào câu lệnh nhắc. Kết quả này mang lại một kết luận thực nghiệm quan trọng: \textbf{việc đơn thuần chèn thêm thông tin siêu dữ liệu hình ảnh dạng văn bản vào ngữ cảnh không giúp mô hình ngôn ngữ nâng cao chất lượng câu trả lời, thậm chí có thể gây phân tán sự tập trung chú ý của mô hình}. Do đó, hệ thống lựa chọn duy trì tắt tùy chọn ngữ cảnh đa phương thức dạng văn bản (\texttt{MULTIMODAL\_CONTEXT\_ENABLED = false}) trong cấu hình vận hành mặc định.

\subsection{Các ưu điểm nổi bật của hệ thống}
\begin{itemize}
    \item \textbf{Khả năng kiểm soát hiện tượng ảo giác hiệu quả:} Điểm Độ trung thực đạt $4{,}56/5$ ở nhóm văn bản và $5{,}00/5$ ở nhóm ngoài phạm vi, minh chứng hệ thống luôn bám sát sách giáo khoa chính thống.
    \item \textbf{Cơ chế trích dẫn xác thực và minh bạch tuyệt đối:} Số trang và tên sách trong trích dẫn được trích xuất trực tiếp mang tính xác định từ siêu dữ liệu của phân đoạn tài liệu được truy xuất, triệt tiêu hoàn toàn khả năng mô hình tự ý tạo dựng số trang sai lệch.
    \item \textbf{Phương pháp đánh giá khách quan, khả năng tái lập cao:} Bộ dữ liệu kiểm thử được gắn nhãn nguồn tường minh, cho phép đo lường chính xác các độ đo IR mà không phụ thuộc vào cảm tính hay biến động của mô hình ngôn ngữ.
    \item \textbf{Tối ưu hóa tài nguyên phần cứng:} Toàn bộ quy trình số hóa và lập chỉ mục cho $2.387$ trang sách giáo khoa vận hành mượt mà trên CPU máy trạm thông thường mà không cần GPU tăng tốc.
    \item \textbf{Bóc tách hình ảnh minh bạch và giải thích được:} Phương pháp neo chú thích đảm bảo mỗi khung hình trích xuất đều liên kết chặt chẽ với một đối tượng định danh chính thống trên trang sách.
\end{itemize}

\subsection{Các giới hạn kỹ thuật và định hướng phát triển}
\label{sec:hanche}
Trên tinh thần khoa học nghiêm túc và cầu thị, đề tài nhận diện rõ các hạn chế kỹ thuật hiện tại:

\begin{enumerate}
    \item \textbf{Quy trình kiểm duyệt bộ kiểm thử:} Toàn bộ $240$ câu hỏi kiểm thử sau khi sinh tự động đã được chuyên gia rà soát tổng thể; tuy nhiên việc bổ sung thêm một lượt đánh giá chéo độc lập từ người thứ hai sẽ giúp thiết lập khoảng tin cậy thống kê hoàn thiện hơn.
    \item \textbf{Sự biến thiên nhẹ giữa các mô hình giám khảo:} Việc luân chuyển giữa bốn mô hình trong tập giám khảo độc lập có thể mang lại một biên độ dao động nhỏ về mặt phân phối điểm số.
    \item \textbf{Lượng hóa chuyên sâu mức độ cải thiện công thức sau ghép dòng:} Giải pháp kiến trúc lai phân cấp đã hoàn thành xử lý trên toàn bộ $16.515$ phân đoạn của kho tri thức ($4\,293$ lượt ghép dòng MinerU thành công); tuy nhiên, việc thực hiện thêm một đợt đo đối sánh tỷ lệ lỗi ký tự trên tập mẫu lớn tiếp theo sẽ giúp lượng hóa chuẩn xác hơn mức cải thiện của chỉ số dưới.
    \item \textbf{Độ phủ của trường siêu dữ liệu số hiệu bài học:} Do nguyên tắc bảo toàn tính chân thực và tránh suy đoán khi mục lục bị gián đoạn, trường số hiệu bài học hiện mới bao phủ $29{,}8\%$ số phân đoạn ($4/12$ cuốn sách).
    \item \textbf{Sự đánh đổi về trục tổng hợp phân tích:} Việc ưu tiên tổng hợp độ đo theo loại câu hỏi nhằm phục vụ mục tiêu đánh giá bản chất tương tác dẫn đến việc lược bỏ bảng so sánh riêng lẻ giữa 12 cuốn sách.
    \item \textbf{Thách thức phân tách các hình ảnh ghép phức hợp:} Một số hình ảnh ghép nhiều tấm liền kề đôi khi vẫn bị trích xuất gộp chung thành một khung hình duy nhất.
    \item \textbf{Cờ cảnh báo cần rà soát tự động:} Do thiết lập điều kiện nghiêm ngặt, cờ cảnh báo rà soát tự động bật ở $70{,}8\%$ số đoạn văn bản, làm giảm khả năng phân loại chọn lọc cho người kiểm duyệt.
    \item \textbf{Hiện tượng nghịch đảo độ phủ tại R@10:} Nguyên nhân khiến chỉ số R@10 của cấu hình đề xuất thấp hơn nhẹ so với cấu hình chưa tái xếp hạng (nghi ngờ do ngưỡng điểm sàn của mô hình tương tác chéo) cần tiếp tục được khảo nghiệm qua các bước quét ngưỡng tham số chi tiết.
    \item \textbf{Bảo toàn cấu trúc hàng và cột của bảng số liệu:} Quá trình OCR phân mảnh hiện tại đôi khi làm xáo trộn thứ tự quan hệ giữa các hàng và cột đối với các bảng thống kê phức tạp.
    \item \textbf{Quy mô thử nghiệm người dùng thực tế:} Các đánh giá hiện tại tập trung vào kiểm thử tự động định lượng; việc triển khai thử nghiệm thực tế rộng rãi trên học sinh và giáo viên trong các lớp học chính quy là nhiệm vụ cần tiếp tục thực hiện.
\end{enumerate}

\chapter{KẾT LUẬN VÀ HƯỚNG PHÁT TRIỂN}
\label{chapter:ket_luan}

\section{Kết luận}
Đề tài tốt nghiệp ``Nghiên cứu phương pháp RAG đa phương thức cho trợ lý ảo thông minh hỗ trợ dạy và học môn Khoa học tự nhiên cấp Trung học cơ sở'' đã được triển khai một cách hệ thống và toàn diện: từ thiết lập cơ sở lý luận, thiết kế kiến trúc phân tầng, chuẩn hóa quy trình trích xuất học liệu đa phương thức, đến tiến hành thực nghiệm đối chiếu định lượng và phân tích chuyên sâu. Đối chiếu với \textbf{năm mục tiêu nghiên cứu} được xác lập tại Chương~\ref{chapter:tong_quan_de_tai}, các kết quả trọng tâm đạt được như sau:

\begin{enumerate}
    \item \textbf{Mục tiêu 1 về xây dựng cơ chế truy vấn đa định dạng:} Đạt được phần lớn các yêu cầu đề ra. Hệ thống đã hiện thực hóa thành công khả năng tiếp nhận và truy xuất đồng thời dữ liệu văn bản lý thuyết cùng tài nguyên hình ảnh với $3.881$ khung hình được bóc tách minh bạch theo phương pháp neo chú thích chính quy. Đối với bài toán đặc thù về công thức định lượng thuộc các phân môn Vật lí và Hoá học, đề tài đã nhận diện rõ điểm nghẽn của công cụ nhận dạng ký tự quang học truyền thống, đồng thời đề xuất và triển khai thành công kiến trúc xử lý lai phân cấp theo dòng trên toàn bộ kho tri thức ($4.293$ lượt ghép dòng MinerU thành công trên $16.515$ đoạn văn bản). Song song với đó, cơ chế chuẩn hoá ký hiệu công thức ở cả phía truy vấn và chỉ mục từ khóa thưa đã nâng tỷ lệ tìm kiếm thành công phân đoạn tài liệu đích từ $1/12$ lên $11/12$ trên tập mẫu kiểm thử.
    
    \item \textbf{Mục tiêu 2 về thiết lập cơ sở tri thức véc-tơ toàn diện:} Hoàn thành xuất sắc. Toàn bộ chương trình môn Khoa học tự nhiên từ lớp 6 đến lớp 9 của cả ba bộ sách giáo khoa chuẩn mực hiện hành, bao gồm 12 cuốn sách tích hợp đầy đủ ba phân môn Vật lí, Hoá học và Sinh học, đã được số hoá và lập chỉ mục trọn vẹn: \textbf{2.387/2.399 trang tài liệu} nội dung bài học được xử lý thành \textbf{20.396 véc-tơ} ($16.515$ véc-tơ văn bản kích thước $1024$ chiều và $3.881$ véc-tơ hình ảnh kích thước $512$ chiều) lưu trữ trong ChromaDB, đồng bộ cùng chỉ mục thưa Okapi BM25 với $20.914$ từ vựng. Toàn bộ quy trình trích xuất, chuyển đổi và nạp dữ liệu này vận hành hiệu quả trực tiếp trên CPU của máy trạm thông thường mà không đòi hỏi tài nguyên GPU.

    \item \textbf{Mục tiêu 3 về tối ưu hóa pha truy xuất tri thức:} Hoàn thành xuất sắc và tạo nên đóng góp phương pháp luận cốt lõi. Kiểm nghiệm định lượng trên bộ dữ liệu kiểm thử $240$ câu hỏi lấy mẫu ngẫu nhiên ở quy mô vận hành thực tế khẳng định cấu hình đề xuất, gồm truy xuất lai kết hợp tái xếp hạng bằng mô hình tương tác chéo, đạt chỉ số \textbf{MRR đạt $0{,}7789$}, vượt trội so với các cấu hình đối chứng: chỉ mục thưa BM25 thuần đạt $0{,}6636$, kênh ngữ nghĩa thuần đạt $0{,}5664$ và phương pháp lai khi chưa tái xếp hạng đạt $0{,}6947$. Các thực nghiệm bóc tách thành phần cũng khẳng định vai trò tích cực của bước tái xếp hạng, đặc biệt nâng chỉ số R@1 của kênh véc-tơ đặc thêm $+53{,}6\%$, đồng thời xác nhận tác động bất lợi của cổng lọc khoảng cách tương đối ở quy mô thực tế, cung cấp cơ sở thực nghiệm vững chắc cho việc thiết lập cấu hình mặc định của hệ thống.

    \item \textbf{Mục tiêu 4 về thiết lập khung đánh giá thực nghiệm đối chiếu:} Hoàn thành toàn diện. Đề tài đã xây dựng thành công bộ dữ liệu kiểm thử $240$ câu hỏi có gắn nhãn nguồn xác thực và tiến hành đánh giá thực nghiệm đa chiều:
    \begin{itemize}
        \item Đối chiếu chi tiết giữa truy xuất lai và BM25 thuần cùng các cấu hình thành phần;
        \item Thực nghiệm đối sánh chuyên biệt giữa RAG đa phương thức và RAG thuần văn bản trên $52$ câu hỏi hình ảnh, mang lại kết luận định lượng quan trọng về việc duy trì cấu hình tối ưu;
        \item Kiểm chứng định lượng năng lực từ chối trả lời chính xác đối với các câu hỏi ngoài phạm vi môn học ($29/30$ câu, đạt tỷ lệ $96{,}67\%$, với điểm số giám khảo về Tính đúng đạt $4{,}87/5$), minh chứng cho khả năng kiểm soát ảo giác hiệu quả của hệ thống.
    \end{itemize}

    \item \textbf{Mục tiêu 5 về hiện thực hóa ứng dụng trợ lý học tập trực quan:} Hoàn thành đầy đủ. Ứng dụng web tương tác trên nền tảng Next.js đã được phát triển hoàn thiện, hỗ trợ người học gửi câu hỏi và nhận câu trả lời dưới dạng dòng dữ liệu phản hồi trực tiếp liên tục, hiển thị đồng bộ hình ảnh minh họa từ máy chủ, kết xuất \textbf{công thức khoa học chuẩn ký hiệu KaTeX} và cung cấp \textbf{khối trích dẫn minh bạch} phân tách rõ tên sách cùng số trang in thực tế.
\end{enumerate}

Về chất lượng câu trả lời sinh ra, đánh giá độc lập theo phương pháp mô hình ngôn ngữ lớn làm giám khảo với hội đồng bốn mô hình luân phiên ghi nhận điểm số trung bình toàn hệ thống đạt mức cao: Tính đúng đạt $4{,}033/5$, Độ trung thực đạt $4{,}417/5$ và Độ liên quan đạt $4{,}492/5$. Trong đó, nhóm câu hỏi ngoài phạm vi đạt điểm số tuyệt đối $5{,}00/5$ về Độ trung thực và Độ liên quan cùng điểm Tính đúng $4{,}87/5$; nhóm câu hỏi văn bản đạt Tính đúng $4{,}19/5$, Độ trung thực $4{,}56/5$ và Độ liên quan $4{,}59/5$. Đặc biệt, cơ chế xây dựng trích dẫn xác thực dựa trên siêu dữ liệu phân đoạn văn bản truy xuất đã triệt tiêu hoàn toàn hiện tượng mô hình ngôn ngữ tự tạo dựng số trang sai lệch, đảm bảo độ tin cậy khoa học vững chắc cho người học và giáo viên khi tra cứu học liệu.

\section{Hạn chế của đề tài}
Bên cạnh những kết quả đạt được, đề tài nhận diện các hạn chế kỹ thuật hiện tại cần được tiếp tục hoàn thiện (chi tiết số liệu tại Mục~\ref{sec:hanche}):
\begin{itemize}
    \item Quy trình kiểm định bộ dữ liệu kiểm thử $240$ câu hỏi hiện dựa trên một lượt rà soát chuyên gia tổng thể; việc bổ sung thêm quy trình thẩm định chéo độc lập từ nhiều chuyên gia chuyên ngành sẽ nâng cao hơn nữa độ chuẩn xác thống kê.
    \item Việc luân chuyển giữa bốn mô hình trong hội đồng giám khảo độc lập có thể tạo ra một biên độ dao động nhỏ về mặt phân bố điểm số.
    \item Quá trình ghép dòng công thức theo kiến trúc lai đã hoàn tất xử lý trên toàn bộ kho tri thức, song việc tiếp tục tiến hành các đợt đo tỷ lệ lỗi ký tự trên tập mẫu mở rộng sẽ giúp lượng hóa sâu sắc hơn mức độ chuẩn hóa của chỉ số dưới.
    \item Trường siêu dữ liệu số hiệu bài học hiện mới bao phủ $29{,}8\%$ số phân đoạn do hệ thống tuân thủ nguyên tắc tránh suy đoán khi mục lục bị gián đoạn bố cục.
    \item Hiện tượng suy giảm nhẹ tại chỉ số R@10 của cấu hình tái xếp hạng cần được tiếp tục phân tích qua các thực nghiệm quét ngưỡng điểm số tối thiểu.
    \item Cấu trúc hàng và cột của một số bảng biểu số liệu phức tạp đôi khi bị xáo trộn qua bước phân mảnh văn bản thông thường.
    \item Hệ thống mới được thử nghiệm chủ đạo qua các kịch bản kiểm thử định lượng tự động; quy mô triển khai thử nghiệm thực tế với giáo viên và học sinh tại các trường phổ thông còn giới hạn.
\end{itemize}

\section{Hướng phát triển tương lai}
Dựa trên những phát hiện khoa học và các mặt hạn chế được chỉ ra, các hướng nghiên cứu và phát triển tiếp theo của đề tài được định hình theo thứ tự ưu tiên như sau:

\begin{enumerate}
    \item \textbf{Triển khai thử nghiệm thực tế trong môi trường sư phạm:} Đưa hệ thống vào thử nghiệm diện rộng tại các cơ sở giáo dục Trung học cơ sở nhằm thu thập phản hồi định tính và định lượng từ đội ngũ giáo viên giảng dạy và học sinh, đánh giá tác động thực tiễn của trợ lý ảo đối với hiệu quả dạy và học môn Khoa học tự nhiên.
    \item \textbf{Đo lường và lượng hoá chuyên sâu chất lượng nhận dạng công thức:} Tiến hành các đợt thực nghiệm đối sánh mở rộng trên tập mẫu lớn văn bản sau ghép dòng để đánh giá định lượng mức cải thiện của chỉ số dưới và độ chuẩn xác của các phương trình hoá học, vật lí.
    \item \textbf{Nghiên cứu tối ưu hóa ngưỡng tham số của bộ tái xếp hạng:} Tiến hành quét lưới chi tiết đối với tham số \texttt{RERANK\_SCORE\_MIN} nhằm giải quyết triệt để hiện tượng nghịch đảo độ phủ tại chỉ số R@10, cân bằng tối ưu giữa độ chính xác và độ phủ của các phân đoạn tài liệu trích xuất.
    \item \textbf{Hoàn thiện trục phân tích mục lục bài học:} Nâng cấp bộ phân tích cấu trúc mục lục cho các bộ sách còn lại để hoàn thiện trường siêu dữ liệu số hiệu bài học đạt độ phủ $100\%$ trên toàn bộ kho tri thức, mở khoá tính năng tra cứu học liệu chuyên biệt theo từng bài học cụ thể.
    \item \textbf{Tích hợp kỹ thuật bảo toàn cấu trúc bảng biểu:} Ứng dụng các giải pháp nhận dạng và phân tích cấu trúc bảng biểu chuyên sâu nhằm bảo toàn nguyên vẹn mối quan hệ hàng và cột của dữ liệu số liệu khi phân mảnh văn bản.
    \item \textbf{Tích hợp mô hình thị giác -- ngôn ngữ trực tiếp:} Nghiên cứu kết hợp các mô hình thị giác -- ngôn ngữ đa phương thức để tiếp nhận trực tiếp hình ảnh sơ đồ và thí nghiệm phức tạp trong câu lệnh nhắc, khắc phục giới hạn của việc chỉ biểu diễn hình ảnh dưới dạng chuỗi siêu dữ liệu văn bản.
    \item \textbf{Lượng tử hóa và tối ưu hóa suy luận mô hình sinh:} Nghiên cứu ứng dụng các kỹ thuật lượng tử hóa trọng số nhằm triển khai mô hình ngôn ngữ lớn cục bộ với chi phí phần cứng tối thiểu, nâng cao tốc độ phản hồi và bảo mật dữ liệu học tập.
\end{enumerate}


\newpage
\phantomsection
\addcontentsline{toc}{chapter}{\textbf{TÀI LIỆU THAM KHẢO}}
\printbibliography[title={TÀI LIỆU THAM KHẢO}]

\end{document}
